\chapter{Notions of commutative algebra}

\section{Modules and rings}
\subsection{Basic definitions}
\begin{dang}
    \begin{definition}\cite{curtis1966representation}
        Let $R$ be a ring and $M$ an abelian group written additively. Then $M$ is said to be an $R-$module if there is a map $R\times M\rightarrow M :(r,m)\mapsto rm$ satisfying:
        \begin{enumerate}
            \item[(1)] $r(sm)=(rs)m$
            \item[(2)] $(r+s)m = rm + sm$
            \item[(3)] $r(m+n) = rm+rn$
            \item[(4)] $1m =m$
        \end{enumerate}
        for any $r,s\in R$ and $m,n\in M$.
    \end{definition}
\end{dang}
This is the notion of left $R-$module. Right $R-$modules are defined in a similar way. If $R$ is a commutative ring, these two notions are equivalent. In this section, modules are all ``left'' (see remarks after definition~(\ref{Definition: definition 1.1.15})).
\begin{dang}
    \begin{definition}\cite{curtis1966representation}
        Given two $R-$modules $M,N$, a \textit{morphism} $f: M\rightarrow N$ from $M$ to $N$, is a map satisfying $f(m+n)=f(m)+f(n)$ and $rf(m)=f(rm)$ for all $m,n\in M$ and $r\in R$. The set of morphisms from $M$ to $N$ is written as $\text{Hom}_R(M,N)$. 
    \end{definition}
\end{dang}

\begin{dang}
    \begin{notation}\cite{jacobsonbasic}
        With these two definitions we have the notion of \textit{category of }$R-$modules, which will be denoted by $\cat{R}$.
    \end{notation}
\end{dang}

\begin{dang}
    \begin{definition}\cite{jacobsonbasic}
         Given a family $\{M_\alpha\}_{\alpha\in A}$ of $R-$modules, the product $M=\prod_{\alpha\in A}M_\alpha$ is defined to be the set of maps $f: A\rightarrow\bigcup_{\alpha\in A}M_\alpha$ with the property $f(\alpha)\in M_\alpha,\;\forall\alpha\in A$. So, $\prod_{\alpha\in A}M_\alpha$ inherits a structure of a $R-$module by defining $rf$ to be the map $\alpha\mapsto rf(\alpha)$. The product is equipped with canonical projections $p_\alpha: M\rightarrow M_\alpha: f\mapsto f(\alpha)$.
    \end{definition}
\end{dang}
The product $M=\prod_{\alpha\in A}M_\alpha$ is characterized the following universal property: Suppose that we are given a $R-$module $N$ with a collection of homomorphisms $g_\alpha: N\rightarrow M_\alpha$, then there exist a unique homomorphism $g: N\rightarrow M$ such that the diagram
\begin{center}
    \begin{tikzcd}
        N \arrow[r, "g"]\arrow[rd, "g_\alpha"]& M\arrow[d,"p_\alpha"]\\
         & M_\alpha
    \end{tikzcd}
\end{center}
commutes for all $\alpha\in A$ \cite{jacobsonbasic}.

\begin{dang}
    \begin{definition} \cite{jacobsonbasic}  
        Given a family $\{M_\alpha\}_{\alpha\in A}$ of $R-$modules, the direct sum $M=\bigoplus_{\alpha\in A} M_\alpha$ is defined to be the submodule of $\prod_{\alpha\in A}M_\alpha$ consisting of maps $f: A\rightarrow \bigcup_{\alpha\in A}M_\alpha$ with finite support, i.e. $f(\alpha)=0$ for all but a finite number of $\alpha \in A$. The direct sum $M$ is equipped with canonical injections $i_\alpha: M_\alpha\rightarrow M$, where $i_\alpha(m_\alpha)\in M$ satisfies $i_\alpha(m_\alpha)(\beta)=0,\;\beta\neq\alpha$ while $i_\alpha(m_\alpha)(\alpha)=m_\alpha$.
    \end{definition}
\end{dang}
The direct sum $M=\bigoplus_{\alpha\in A}M_\alpha$ satisfies the following universal property: Given a $R-$module $N$ with a family of maps $f_\alpha: M_\alpha\rightarrow N$, then there exists an unique homomorphism $f: M\rightarrow N$ such that the diagram
\begin{center}
    \begin{tikzcd}
        N & M\arrow[l, "f" ]\\
           & M_\alpha\arrow[ul, "f_\alpha"]\arrow[u, "i_\alpha"]
    \end{tikzcd}
\end{center}
commutes for all $\alpha\in A$ \cite{jacobsonbasic}. 


\begin{dang}
    \begin{definition}\cite{curtis1966representation}
        Let $S$ be a subset of a $R-$module $M$. Then $S$ is said to be $R-$\textit{free} if whenever $m_1,\dots m_k\in S$ and
        \[
        r_1m_1+\dots+r_km_k=0 
        \]
        for some $r_1,\dots,r_k\in R$. Then $r_i=0\quad\forall i=1,\dots,k$. 
    \end{definition}
\end{dang}
\begin{dang}
    \begin{definition}\cite{curtis1966representation}
        A subset $S$ of an $R-$module $M$ is said to be a set of generators of $M$, if for every $m\in M$, we have $r_1,\dots,r_k\in R$ and $m_1,\dots,m_k\in S$ such that $m=\sum_{i=1}^kr_im_i$. 
    \end{definition}
\end{dang}
\begin{dang}
    \begin{notation}
        Let $M$ be an $R-$module with a finite set of generators $\{m_1,\dots,m_k\}$. In this case we will write $M = \sum_{i=1}^kRm_i$.
    \end{notation}
\end{dang}
\begin{dang}
    \begin{definition}\cite{curtis1966representation}
        An $R-$module $M$ is said to be \textit{cyclic} if it is generated by a unique element $m\in M$. In this case we can just write $M=Rm$.
    \end{definition}
\end{dang}
\begin{dang}
    \begin{definition}\cite{curtis1966representation}
        An $R-$module $M$ is said to be a \textit{free} $R-$module if it has a $R-$free set of generators $S$. We will say that $S$ is an $R-$basis of $M$ and that $M$ is \textit{freely generated} on $S$.
    \end{definition}
\end{dang}
\begin{dang}
    \begin{corollary}\cite{curtis1966representation}
       Let $M=\sum_{i=1}^kRm_i$ be a free $R-$module. Then 
       \[
       M\cong \bigoplus_{i=1}^{k}Rm_i\cong\bigoplus_{i=1}^{k}R=R^k
       \]
    \end{corollary}
\end{dang}
Proof: Define the surjective homomorphism $R^k\rightarrow M:(r_1,\dots,r_k)\mapsto r_1m_1+\dots+r_km_k$. Since $\{m_1,\dots,m_k\}$ is free, this map is injective and we have an isomorphism.

\begin{dang}
    \begin{definition}
    Let $(A, \leq)$ be a partially ordered set, $\{M_\alpha\}_{\alpha\in A}$ a family of $R-$modules 
    and $\{f_{\alpha\beta}\}_{\alpha\leq\beta}$ a family of homomorphisms $f_{\alpha\beta}:M_\alpha\rightarrow M_\beta$ 
    satisfying $f_{\alpha\alpha}=\mathds{1}_{M_\alpha}\;\forall\alpha\in A$ 
    and $f_{\alpha\gamma}=f_{\beta\gamma}\circ f_{\alpha\beta}$ for $\alpha\leq\beta\leq\gamma$. 
    Then we say that the pair $(\{M_{\alpha}\}_{\alpha\in A},\{f_{\alpha\beta}\})$ 
    is a \textit{directed system}. The direct limit of a directed system 
    $\{(M_{\alpha}\}_{\alpha\in A},\{f_{\alpha\beta}\})$ is defined to be the module
    \begin{equation}
        \varinjlim M_\alpha = \bigcup_{\alpha\in A}M_\alpha\bigg/\sim
    \end{equation}
    where $m_\alpha\in M_\alpha\sim m_\beta\in M_\beta$ if there exists $\gamma\geq \alpha,\beta$ such that $f_{\alpha\gamma}(m_\alpha)=f_{\beta\gamma}(m_\beta)$. 
    \end{definition}
\end{dang}
Let $\alpha,\beta \in A$ and take $\gamma\geq\alpha,\beta$. Then, we define the sum
\begin{equation}
    [m_\alpha]+[m_\beta]=[f_{\alpha\gamma}(m_\alpha)+f_{\beta\gamma}(m_\beta)]
\end{equation}
Moreover, the $R-$action on $\varinjlim M_\alpha$ is defined as $r[m_\alpha]=[rm_\alpha]$. All these operations are 
clearly well defined.

\subsection{Tensor product of modules}
The main purpose of this subsection is to construct the tensor product modules. 
\begin{dang}
    \begin{definition}\cite{curtis1966representation}
        Let $M$ (resp. $N$) be right (resp. left) $R-$modules and $P$ an abelian group. A map $f:M\times N\rightarrow P$ is said to be \textit{balanced} if it satisfies
        \begin{enumerate}
            \item[(1)] $f(m_1+m_2,n)=f(m_1,n)+f(m_2,n)$
            \item[(2)] $f(m,n_1+n_2)=f(m,n_1)+f(m,n_2)$
            \item[(3)] $f(mr,n)=f(m,rn)$
        \end{enumerate}
        for all $r\in R$ and $m_i\in M,\;n_i\in N$.
    \end{definition}
\end{dang}
This concept can be seen as an extension of the concept of bilinear maps from linear algebra. 
\begin{dang}
    \begin{theorem}\label{theorem: existence of tensor products}\cite{curtis1966representation}
        Let $M$ (resp. $N$) be right (resp. left) $R-$modules. Then there exist an abelian group $M\otimes_RN$, called the \textit{tensor product} between $M$ and $N$, with the following properties. 
        \begin{enumerate}
            \item[(1)] There exist a balanced map $\pi: M\times N\rightarrow M\otimes_R N$ such that every element of $M\otimes_R N$ is of the form $\sum_{i=1}^k\pi(m_i,n_i)$ where $m_i\in M$ and $n_i\in N$.
            \item[(2)] Every balanced map $f: M\times N\rightarrow P$ factors through $M\otimes_R N$, i.e. there exist $\Tilde{f}: M\otimes_RN\rightarrow P$ such that
            $\Tilde{f}\circ\pi=f$.
        \end{enumerate}
    \end{theorem}   
\end{dang}
Proof: Let $A=\bigoplus_{M\times N}\bb{Z}$, i.e. $A$ is the free $\bb{Z}-$module generated by $M\times N$. Writing $\overline{(m,n)}\in A$ in place of $i_{(m,n)}(m,n)$,  let $G\subset A$ be the submodule generated by elements of the form:
\begin{equation}
\left\{\begin{matrix}
    \overline{(m_1+m_2,n)}-\overline{(m_1,n)}-\overline{(m_2,n)}\\
    \overline{(m,n_1+n_2)}-\overline{(m,n_1)}-\overline{(m,n_2)}\\
    \overline{(mr,n)}-\overline{(m,rn)}
\end{matrix}\right.
\end{equation}
So we define the tensor product to be the $\bb{Z}-$module $M\otimes_R N \bydef A/G$. By definition, it is clear that $\pi:M\times N\rightarrow M\otimes_R N:(m,n)\mapsto \overline{(m,n)}+G$ is balanced, and that $M\otimes_R N$ is generated by sums of the form $\sum_{i=1}^k\pi(m_i,n_i)$ where $m_i\in M$. Now, let $f:M\times N\rightarrow P$ be a balanced map. Since $A$ is the free $\bb{Z}-$module generated by $M\times N$, we have a map $\bar{f}: A\rightarrow P:\overline{(m,n)}\mapsto f(m,n)$. But since $f$ is balanced $G\subset \text{Ker}\bar{f}$, so by the isomorphism theorem we get a $\bb{Z}-$module homomorphism $\Tilde{f}: A/G=M\otimes_R N\rightarrow P: \pi(m,n)\mapsto f(m,n)$.
\begin{dang}
\begin{notation}
    From now on, in the conditions of the previous theorem, we will write $m\otimes n$ instead of $\pi(m,n)$.
\end{notation}
\end{dang}
It's important to note the properties (1) and (2) in theorem~(\ref{theorem: existence of tensor products}) characterize the tensor product uniquely. In fact, suppose that $M\otimes_R^\prime N$ is another abelian group that satisfies both properties (1) and (2). Let $\pi^\prime:M\times N\rightarrow M\otimes_R^\prime N$ be the balanced map such that $M\otimes_R^\prime N$ is generated by elements of the form $\sum_{i=1}^k\pi^\prime(m_i,n_i)$ for $m_i\in M$ and $n_i\in N$. So we have the induced homomorphisms $\overline{\pi}: M\otimes_RN\rightarrow M\otimes_R^\prime N$ and $\overline{\pi}^\prime:M\otimes_R^\prime N\rightarrow M\otimes_RN$ such that 
\begin{equation}
    \overline{\pi}^\prime(\pi(m,n))=\pi^\prime(m,n)\;\;\text{ and }\;\;\overline{\pi}(\pi^\prime(m,n))=\pi(m,n)
\end{equation}
By the property (2), both maps $\overline{\pi}$ and $\overline{\pi}^\prime$ are surjective and therefore, inverses of each other by the last equation. Particularly, this proves that $\overline{\pi}: M\otimes_RN\rightarrow M\otimes_R^\prime N$ is an isomorphism of groups.
\begin{dang}
    \begin{definition}\cite{curtis1966representation}\label{Definition: definition 1.1.15}
        Let $S$ and $R$ be rings. A $(S,R)-$bimodule $M$ is a module that is both a left $S-$module and a right $R-$module that satisfies the additional property 
        \begin{equation}
            (sm)r = s(mr)\quad\forall s\in S, r\in R \text{ (``associativity of actions'')}
        \end{equation}
    \end{definition}
\end{dang}
An important observation is that if $R$ is a commutative ring, then there is no difference between left $R-$modules and right $R-$modules. In fact, any left (or right) $R-$module is a $(R,R)-$bimodule. For example, if $M$ is a left $R-$module, we can just define $mr\bydef rm$ turning $M$ into a $(R,R)-$bimodule. The associativity property is a direct consequence of the commutativity of $R$. Therefore, when working with modules over commutative rings we will not specify left or right. For example, any vector space $V$ over a field $\bb{F}$ is a $(\bb{F},\bb{F})-$bimodule and if $W$ is another vector space over the same field, we conclude that $V\otimes_{\bb{F}} W$ is also a vector space over $\bb{F}$.\par
Now, let $M$ be an $(S,R)-$bimodule and $N$ a right $R-$module. We claim that $M\otimes_RN$ has a structure of a left $S-$module. To see this, just define for $s\in S$ the map $f_s: M\times N\rightarrow M\otimes_RN:(m,n)\mapsto sm\otimes n$. This map is clearly balanced, so this define an action $s(m\otimes n)\bydef \Tilde{f_s}(m\otimes n)=(sm)\otimes n$. 
In summary, we have:
\begin{dang}
    \begin{theorem}\cite{curtis1966representation}
        Let $M$ be a $(S,R)-$bimodule and $N$ a right $R-$module. The tensor product of modules $M\otimes_R N$ is a left $S-$module with action
        \begin{equation}
            s(m\otimes n)=(sm)\otimes n.
        \end{equation}
    \end{theorem}
\end{dang}
 Similarly, if $N$ is a left $S-$module, then $N\otimes_S M$ is a right $R-$module with action $(n\otimes m)r=n\otimes(mr)$. The proof is exactly the same as for the module $M\otimes_R N$.\par
Now we come to important application of the definition of tensor product. If $R$ is a ring contained in another ring $S$, then we can define a functor $\text{Ind}_R^S:\cat{R}\rightarrow\cat{S}$ called \textit{induction}. The functor is defined as follows, for each left $R-$module $M$, we let  $\text{Ind}_R^SM\bydef S\otimes_RM$. Since $S$ is clearly a $(S,R)-$bimodule, the abelian group $\text{Ind}_R^SM$ is also a left $S-$module. Also, for $f\in\text{Hom}_R(M,N)$ we have 
\begin{equation}
    \text{Ind}_R^Sf:\text{Ind}_R^SM\rightarrow\text{Ind}_R^SN:s\otimes m\mapsto s\otimes f(m).
\end{equation}
An important property of induction is that it is transitive.
\begin{dang}
    \begin{theorem}\label{theorem: associativity of the induction functor}
        Let $R\subset S\subset T$ be rings and $M$ a left $R-$module. Then we have an isomorphism of left $T-$modules
        \begin{equation}
            \text{Ind}_S^T\text{Ind}_R^SM\cong\text{Ind}_R^TM
        \end{equation}
    \end{theorem}
\end{dang}
To show this, we will need some lemmas.
\begin{dang}
    \begin{lemma}\label{lemma: M otimes N isomorphic to M otimes N'}
        Let $N\cong N^\prime$ be two isomorphic left $R-$modules and $M$ a $(S,R)-$bimodule. Then we have an isomorphism of left $S-$modules $M\otimes_RN\cong M\otimes_RN^\prime$.
    \end{lemma}
\end{dang}
Proof: Let $f:N\rightarrow N^\prime$ be an isomorphism of left $R-$modules and $g$ its inverse. Then $M\times N\rightarrow M\otimes_R N^\prime:(m,n)\mapsto m\otimes f(n)$ is balanced, so we have an induced homomorphism of groups $\varphi:M\otimes_R N\rightarrow M\otimes_R N^\prime:m\otimes n\mapsto m\otimes f(n)$. It's clear that $\varphi$ is also a homomorphism of modules. By the same reasoning, we can construct a homomorphism $\psi: m\otimes n^\prime\mapsto m\otimes g(n^\prime)$ from $M\otimes_RN^\prime$ to $M\otimes_RN$. But $\varphi$ and $\psi$ are clearly inverses of each other, so we get that $\varphi$ is an isomorphism and we have proved the theorem.
\begin{dang}
    \begin{lemma}\label{lemma: R otimes S is isomorphic to S}
        Let $M$ be a left $R-$module. Since $R$ is a $(R,R)-$bimodule, we can form the tensor product $R\otimes_R M$. We claim that $R\otimes_RM\cong M$ as left $R-$modules.
    \end{lemma}
\end{dang}
Proof: Define the map $R\times M\rightarrow M:(r,m)\mapsto rm$. This map is clearly balanced so we have a homomorphism of abelian groups $\varphi:R\otimes_RM\rightarrow M:r\otimes m\mapsto rm$. Conversely, the map $\psi:S\rightarrow R\otimes_R M:m\mapsto 1\otimes m$ is clearly a group homomorphism. Also we have that $\psi\circ \varphi=1_{R\otimes_R M}$ because 
\begin{equation}
    \psi\circ \varphi(r\otimes m)=\psi(rm)=1\otimes rm=r\otimes m.
\end{equation}
Moreover
\begin{equation}
    \varphi\circ\psi(m)=\varphi(1\otimes m)=1m=m
\end{equation}
so $\psi\circ \varphi=1_{R\otimes_R M}$ and $\varphi\circ\psi=1_{M}$. This proves that $R\otimes_R M$ and $M$ are isomorphic as abelian groups. But both $\varphi$ and $\psi$ are left $R-$modules homomorphisms. So we get that we have in fact, $\varphi$ is an isomorphism of left $R-$modules.
\begin{dang}
    \begin{lemma}\label{lemma: associativity of tensor product of modules}
        Let $M$ be a right $R-$module, $N$ a $(R,S)-$bimodule and $L$ a left $S-$module. Then we have an isomorphism of abelian groups $M\otimes_R(N\otimes_S L)\cong (M\otimes_RN)\otimes_SL$.
    \end{lemma}
\end{dang}
Proof: Fix $m\in M$ and define the map $f_m:N\times L\rightarrow (M\otimes_RN)\otimes_SL:(n,l)\mapsto (m\otimes n)\otimes l$. This map is clearly balanced, so we have a well defined homomorphism of groups 
\begin{equation}
    \Tilde{f_m}:N\otimes_SL\rightarrow(M\otimes_RN)\otimes_SL:n\otimes l\mapsto (m\otimes n)\otimes l.
\end{equation}
This shows that the map
\begin{equation}
    M\times N\otimes_SL\rightarrow (M\otimes_RN)\otimes_SL: (m,n\otimes l)\mapsto \Tilde{f}_m(n\otimes l)
\end{equation}
is well defined. This map is clearly balanced, so we have an induced homomorphism of groups
\begin{equation}
    \varphi:  M\otimes_R( N\otimes_SL)\rightarrow (M\otimes_RN)\otimes_SL: m\otimes(n\otimes l)\mapsto \Tilde{f_m}(n\otimes l)=(m\otimes n)\otimes l
\end{equation}
which is also a left $R-$module homomorphism. Similarly, we can show that 
\begin{equation}
    \psi: (M\otimes_RN)\otimes_SL\rightarrow M\otimes_R( N\otimes_SL):(m\otimes n)\otimes l\mapsto m\otimes(n\otimes l)
\end{equation}
is a well defined left $R-$module homomorphism. Since $\psi$ and $\varphi$ are inverses of each other, we have proved the theorem.\par
Now we can proceed to prove theorem~(\ref{theorem: associativity of the induction functor}). By definition, we just need to check that $T\otimes_S(S\otimes_RM)\cong T\otimes_RM$ as left $T-$modules. But
\begin{equation}
    \begin{split}
        T\otimes_S(S\otimes_RM)&\cong (T\otimes_SS)\otimes_RM\quad\text{(by lemma~(\ref{lemma: associativity of tensor product of modules}))}\\
        &\cong T\otimes_RM\quad\text{(by lemma~(\ref{lemma: M otimes N isomorphic to M otimes N'}) and lemma~(\ref{lemma: associativity of tensor product of modules})})
    \end{split}
\end{equation}\par
Before we finish this subsection, we will need one more useful property of tensor products.
\begin{dang}
    \begin{lemma}
        Let $M_j$ (resp. $N_j$) be right (resp. left) $R-$modules $j=1,2$. If $T_j\in \text{Hom}(M_j,N_j)$, then the map
        \begin{equation}
          T_1\otimes T_2: M_1\otimes_RM_2\rightarrow N_1\otimes_RN_2 :m_1\otimes m_2\mapsto T_1(m_1)\otimes T_2(m_2)
        \end{equation}
        is a well defined homomorphism of groups.
    \end{lemma}
\end{dang}
Proof: We just need to check that the map $(m_1,m_2)\mapsto T_1(m_1)\otimes T_2(m_2)$ is balanced, which is clear. 
\begin{dang}
    \begin{theorem}\label{theorem: tensor product distributes under coproducts}
        Let $M=M_1\oplus M_2$ be a right $R-$module and $N$ a left $R-$module. Then, we have the following group isomorphism $M\otimes_R N\cong M_1\otimes_R N\oplus M_2\otimes_R N$.
    \end{theorem}
\end{dang}
Proof: The direct sum $M=M_1\oplus M_2$ is characterized by the existence of module homomorphisms $i_j:M_j\rightarrow M$ and $p_j:M\rightarrow M_j$ such that 
\begin{equation}
    \left\{\begin{matrix}
        p_j\circ i_j=\mathds{1}_{M_j}\quad j=1,2\\
        \quad p_k\circ i_j=0\quad\text{for }\;\;k\neq j\\
        \sum_{j=1}^2i_j\circ p_j=\mathds{1}_M
    \end{matrix}\right.
\end{equation}
Now, we define $\Tilde{p_j}=i_j\otimes \mathds{1}_N$ and $\Tilde{p_j}=p_j\otimes \mathds{1}_N$. So we get
\begin{equation}
    \begin{split}
        &\Tilde{p_j}\circ\Tilde{i_j}=(p_j\circ i_j)\otimes \mathds{1}_N = \mathds{1}_{M_j}\otimes \mathds{1}_N=1_{M_j\otimes_R N}\quad j=1,2\\
        &\Tilde{p_k}\circ\Tilde{i_j}=(p_k\circ i_j)\otimes \mathds{1}_N =0\otimes \mathds{1}_n=0\quad k\neq j\\
        &\sum_{j=1}^2\Tilde{i_j}\circ\Tilde{p_j}=\left(\sum_{j=1}^2 i_j\circ p_j \right)\otimes \mathds{1}_N=\mathds{1}_M\otimes\mathds{1}_N=\mathds{1}_{M\otimes_RN}
    \end{split}
\end{equation}
Therefore $M\otimes_RN\cong M_1\otimes_RN\oplus M_2\otimes_RN$.
\begin{dang}
    \begin{corollary}\label{corollary: basis of a tensor product of vector spaces}
        Let $V$ and $W$ be two vector spaces of finite dimension over a field $\bb{F}$. Then $\text{dim}_{\bb{F}}(V\otimes_{\bb{F}} W)=\text{dim}_{\bb{F}}V\text{dim}_{\bb{F}}V$. Moreover, if $\{v_i\}$ and $\{w_j\}$ are basis for the vector spaces $V$ and $W$ respectively, then $\{v_i\otimes w_j\}$ is a basis for $V\otimes_{\bb{F}} W$.
    \end{corollary}
\end{dang}
Proof: Let $r=\text{dim}_{\bb{F}}V$, then we get $V\cong \bb{F}^r$ and applying the theorem~(\ref{theorem: tensor product distributes under coproducts}) multiple times, we get the vector space isomorphism isomorphism
\begin{equation}
    V\otimes_{\bb{F}}W \cong \bigoplus_{j=1}^r\bb{F}\otimes_{\bb{F}} W
\end{equation}
But, by the lemma~(\ref{lemma: R otimes S is isomorphic to S}) we have $\bb{F}\otimes_{\bb{F}}W\cong W$. So we have
\begin{equation}
     V\otimes_{\bb{F}}W \cong \bigoplus_{j=1}^r W
\end{equation}
and $\text{dim}_{\bb{F}}(V\otimes_{\bb{F}}W)=\sum_{j=1}^r\text{dim}_{\bb{F}}W=r\text{dim}_{\bb{F}}W=\text{dim}_{\bb{F}}V\text{dim}_{\bb{F}}W$. Moreover, any $v\otimes w\in V\otimes_{\bb{F}}W$ is a sum of the form $\sum_{i,j}\alpha_i\beta_jv_i\otimes w_j$ and $\text{Span}\{v_i\otimes w_j\}= V\otimes_{\bb{F}}W$. But the cardinality of the set $\{v_i\otimes w_j\}$ is exactly the product of dimensions of $V$ and $W$. Therefore, this set must be a basis of $V\otimes_{\bb{F}}W$.
\begin{notation}
    To ease notation, from now on, when working with tensor products between vector spaces over a field, say $\bb{F}$, we will avoid the subscript $\bb{F}$ in the product $\otimes$. 
\end{notation}
\begin{dang}
    \begin{theorem}\label{theorem: tensor of product of algebras is a tensor product}\cite{curtis1966representation}
        Let $A$ and $A^\prime$ be algebras over a field $\bb{F}$. We claim that $A\otimes A^\prime$ is an algebra over $\bb{F}$ with product defined componentwise:
        \begin{equation}
            \left(\sum_ia_i\otimes a_i^\prime\right)\left(\sum_j b_j\otimes b_j^\prime\right)=\sum_{i,j}a_ib_j\otimes a_i^\prime b_j^\prime.
        \end{equation}
    \end{theorem}
\end{dang}
Proof: We will show that if $\sum_{i=1}^na_i\otimes a_i^\prime=0$, then $\sum_{i,j}a_ib_j\otimes a_i^\prime b_j^\prime=0$. This shows that the product is well defined and associativity follows from associativity of $A$ and $A^\prime$. Let $\{v_k\}$ (resp. $\{v_l^\prime\}$)  be a basis of $A$ (resp. $A^\prime$), so $a_i=\sum_k\alpha_{ik}v_k$ and $a_i^\prime=\sum_l\alpha_{il}^\prime v_l^\prime$. Therefore
\begin{equation}
    \begin{split}
        &\sum_{i=1}^n a_i\otimes a_i^\prime=0\Rightarrow\\
        &\quad\sum_{i=1}^n\sum_{k,l}\alpha_{ik}\alpha_{il}^\prime v_{k}\otimes v_l^\prime = 0\Rightarrow\\
        &\quad\;\sum_{i=1}^n\alpha_{ik}\alpha_{il}^\prime=0\;\forall k,l\;\;(\text{because of corollary~(\ref{corollary: basis of a tensor product of vector spaces})})
    \end{split}
\end{equation}
and
\begin{equation}
    \begin{split}
        \sum_{i,j}a_ib_j\otimes a_i^\prime b_j^\prime &= \sum_{i,k,l}\alpha_{ik}\alpha_{il}^\prime v_kb_j\otimes v_l^\prime b_j^\prime\\
        &=\sum_{k,l}\left(\sum_{i}\alpha_{ik}\alpha_{il}^\prime\right) v_kb_j\otimes v_l^\prime b_j^\prime=0.
    \end{split}
\end{equation}
Another important property of tensor products is that its commutes with direct limits.
\begin{dang}
    \begin{theorem}\label{theorem: direct limit commute with tensor product}
        Let $(M_n, f_{mn})$ be a directed system of right $R-$modules and write $M=\varinjlim M_n$ for the associated limit. Then, if $N$ is a left $R-$module, we have an isomorphism $M\otimes_RN\cong \varinjlim M_n\otimes_RN$.
    \end{theorem}
\end{dang}
Proof: Since $(M_n\otimes_RN,f_{mn}\otimes \mathds{1}_N)$ is a directed system, we can form the direct limit $\varinjlim M_n\otimes_R N$. Let $\varphi_n:M_n\rightarrow M:m\mapsto [m]$ be the canonical mappings associated to the direct limit $M$. We will show that $M\otimes_RN\cong \varinjlim M_n\otimes_R N$ by showing that  $(M\otimes_RN, \varphi_n\otimes \mathds{1}_N)$ is an universal ``repelling'' target. Said this, we need to show that if $(Y,g_n)$ is a target, then exists a unique $\alpha$ such that the following diagram commutes:
\begin{center}
    \begin{tikzcd}
        M_m\otimes N \arrow[rr, "f_{mn}\otimes \mathds{1}_N"]\arrow[rd, "\varphi_m\otimes \mathds{1}_N"]\arrow[rdd, bend right,"g_m"] &  &M_n\otimes N\arrow[ld, "\varphi_n\otimes \mathds{1}_N"]\arrow[ldd, bend left, "g_n"]  \\
          & M\otimes_RN \arrow[d, dotted ,"\alpha"] &  \\
          & Y &  
    \end{tikzcd}
\end{center}
Define the map 
\begin{equation}
    M\times N\rightarrow Y: (m,n)\mapsto g_i(m_i\otimes n)
\end{equation}
where $m_i$ is a representative for $m$, i.e $m=[m_i]=\varphi_i(m_i)$. This map is well defined. In fact, let $m_i\sim m_j\Rightarrow f_{li}(m_i)=f_{lj}(m_j)$ for some $l\geq i,j$. By the assumption of the pair $(Y, g_n)$ being a target, we have $g_{m}=g_n\circ (f_{nm}\otimes \mathds{1}_N)$ whenever $n\geq m$. Therefore $g_i(m_i\otimes n)=g_l(f_{li}(m_i)\otimes n)$ and $g_j(m_j\otimes n)=g_l(f_{lj}(m_j)\otimes n)$, but since $f_{il}(m_i)=f_{jl}(m_j)$ we get that $g_i(m_i\otimes n)=g_j(m_j\otimes n)$. This map is clearly balanced, so we get a well defined group homomorphism 
\begin{equation}
    \alpha: M\otimes N\rightarrow Y: m\otimes n\mapsto g_i(m_i\otimes n)
\end{equation}
where $m_i$ is a representative for $m\in M$. For uniqueness, starting with the previous commutative diagram, the equality of compositions $\alpha\circ(\varphi_i\otimes\mathds{1}_N)=g_i$ implies the formula previously defined for $\alpha$. It's clear that if  we have in addition that $(M_n, f_{mn})$ is a directed system of $(S,R)-$modules, then $\alpha$ is a homomorphism of $S-$modules.
\begin{dang}
    \begin{theorem}\label{corollary: tensor distributes over infinite coproducts}
        Let $M=\bigoplus_{\alpha\in A}M_\alpha$ be a direct sum of right $R-$modules. Then, if $N$ is a left $R-$module, we have an isomorphism $M\otimes_R N\cong \bigoplus_{\alpha \in A}M_\alpha\otimes_RN$.
    \end{theorem}
\end{dang}
Proof: Let $\mathcal{A}$ be the set consisting of finite subsets of $A$. For $S\subset A$, write $M_S=\bigoplus_{\alpha\in S}M_\alpha$. Then, we have a canonical monomorphism (in fact an inclusion) $f_{ST}:M_S\rightarrow M_T$ whenever $S\subset T$. Let $M$ be the limit of the directed system $(M_S, f_{ST})$. Now, let $Y$ be a right $R-$module, homomorphisms $g_\alpha: M_\alpha \rightarrow Y\;(\alpha\in A)$ and $i_\alpha\bydef f_{\{\alpha\}A}$. Then we claim that exists a unique homomorphism $g: M\rightarrow Y$ such that $g\circ i_\alpha = g_\alpha$. Let $x\in M$, so $x\in M_S$ for some $S\in\mathcal{A}$. In this case we put $g(x)=g_S(x)$ where $g_S$ is the map induced by the universal property of the direct sum $\bigoplus_{\alpha\in S}M_\alpha$. It's clear that $g$ is well defined and that $g\circ i_\alpha =g_\alpha\;\forall\alpha\in A$. Therefore $M$ satisfies the universal property for the direct sum of the modules $\{M_\alpha\}$, so $M\cong \bigoplus_{\alpha\in A}M_\alpha$. By the theorem~(\ref{theorem: tensor product distributes under coproducts}) we get the result.




\subsection{Structure theorem for finitely generated modules over a principal ideal domain}
\begin{dang}
\begin{definition}\cite{lidl1994introduction}
    A ring $R$ is called an integral domain, if it is commutative and has no divisors of zero. This means that for $a,b\in R$, then if $ab\in R\Rightarrow a\in R$ or $b\in R$.
\end{definition}
\end{dang}
\begin{dang}
    \begin{definition}\cite{lidl1994introduction}
        A commutative ring $R$ such that every nonzero element is a unit is called a \textit{field}.
    \end{definition}
\end{dang}
An important property of integral domains that we will need on the next chapter is the following.
\begin{dang}
    \begin{theorem}\label{a finite integral domain is a field}\cite{lidl1994introduction}
        A finite integral domain is a field.
    \end{theorem}
\end{dang}
Proof: Since $R$ is commutative, because it's a integral domain, we just need to show that every nonzero element in $R$ is a unit. Let $R=\{a_1,\dots,a_n\}$ and $a\in R-\{0\}$. Then, since $R$ is an integral domain, $aa_i\in R$ are all distinct for different $i$. Therefore $aa_i=1$ for some $i$. This shows that every nonzero element $a\in R$ is a unit.
\begin{dang}
    \begin{definition}\cite{hungerford2012algebra}
        A ring $R$ is called principal ideal ring if every ideal $I\subset R$ is of the form $I=(a)$ for some $a\in R$. Ideals of this form will be called \textit{principal ideals}. 
    \end{definition}
\end{dang}
A useful property of principal ideal rings, is that every ascending chain of ideals $I_1\subset I_2\subset \dots$ must stabilize. This means that exists $i\geq1$ such that $I_j=I_i$ whenever $j\geq i$. In fact, $I=\bigcup_i I_i$ is an ideal. Now, since $R$ is principal, we have $I=(a)$. So we see that $a\in I_i$ for some $i$, which implies that $(a)\subset I_i\Rightarrow (a) = I_i$ and $I_j=I_i=(a)$ for all $j\geq i$. Rings with the property that every ascending chain of ideals must terminate are called \textit{noetherian} rings. From the previous discussion we get the following and important result.
\begin{dang}
    \begin{theorem}\label{principal_ideal_is_noetherian}\cite{hungerford2012algebra}
        Every principal ideal ring is noetherian.
    \end{theorem}
\end{dang}
We also have the notion of noetherian modules, that will be useful later.
\begin{dang}
    \begin{definition}\cite{curtis1966representation}
        A module $M$ is said to be \textit{noetherian} if every increasing chain of submodules $M_1\subset M_2\subset \dots$ terminates at some point.
    \end{definition}
\end{dang}
\begin{dang}
    \begin{example}
        Every finite dimensional vector space is noetherian.
    \end{example}
\end{dang}
In this section we will be interested in rings that are both principal ideal rings and integral domains.
\begin{dang}
    \begin{definition}\cite{curtis1966representation}
        A ring $R$ is said to be a principal ideal domain, if it's an integral domain and a principal ideal ring.
    \end{definition}
\end{dang}
\begin{dang}
    \begin{definition}
        Let $R$ be a ring, a proper ideal $I\subset R$ is said to be \textit{maximal} if whenever $I\subset J,\; J$ an ideal, we have necessarily $J=I$ or $J=R$. Also, a proper ideal $P\subset R$ is said to be \textit{prime} if whenever $ab\in P\Rightarrow a\in P$ or $b\in P$.
    \end{definition}
\end{dang}

An important property of principal ideal domains is the following.
\begin{dang}
\begin{theorem}\label{In_PID_every_primeIdeal_is_maximal}
    If $R$ is a principal ideal domain, then $P\subset R$ is a prime ideal if and only if it is a maximal ideal.
\end{theorem}
\end{dang}
Proof: See~(\cite{lidl1994introduction} p.17).
\begin{dang}
    \begin{theorem}\label{submodule_ofa_free_module_is_free}\cite{curtis1966representation}
        Let $R$ be a principal ideal ring and $M$ a finitely generated module over $R$ with basis $\{m_1,\dots,m_k\}$. Then, any submodule $N\subset M$ has a basis of cardinality $s\leq k$.
    \end{theorem}
\end{dang}
Proof: The result is immediate for $k=1$. Indeed $M\cong R$ as $R-$modules, and a submodule of $R$ is the same thing as an ideal. So the principal ideal condition implies the result directly. Now suppose $k>1$ and the result is true for free $R-$modules generated by $k-1$ elements. Since $M=\bigoplus_{i=1}^{k}Rm_i$, we have the projection homomorphism $p_1:M\rightarrow R$. Since $N\subset M$ is a submodule, $p_1(N)\subset R$ is an ideal. But $R$ is a principal ideal ring, so $p_1(N)=(\sigma)$ and there exists $n_1\in N$ such that $p_1(n_1)=\sigma$. Now define $M^\prime=\bigoplus_{i=2}^{k}Rm_i$. Then $p_1(M^\prime)=0\Rightarrow Rn_1\cap M^\prime=(0)$. Also, if $n\in N$ and $p_1(n)=\alpha\sigma$, then $p_1(n-\alpha n_1)=0$, so $n-\alpha n_1\in N\cap M^\prime$. Therefore $N=Rn_1\oplus N\cap M^\prime$. But, by the induction hypothesis, $N\cap M^\prime$ has a basis of cardinality at most $k-1$ elements and the result follows.
\begin{dang}
    \begin{theorem}\cite{curtis1966representation}
        Let $M$ be a free $R-$module finitely generated by a basis $\{m_1,\dots,m_k\}$. The number $k$ is well defined and is called the \textit{rank} of $M$.
    \end{theorem}
\end{dang}
Proof: We have $M\cong\bigoplus_{i=1}^kR$. Let $I\subset R$ be a maximal ideal, then $IM\subset M$ is a submodule isomorphic to $\bigoplus_{i=0}^kI$. Therefore $M/IM\cong \bigoplus_{i=1}^kR/I \cong L^k$ where $L=R/I$ is a field. So $k=\text{dim}_L(M/IM)$ is a well defined number.

\begin{dang}
    \begin{notation}
       Given a ring $R$, we will write $M_{mn}(R)$ for the $m\times n$ matrices over $R$ and $M_n(R)$ for $n\times n$ matrices over $R$.
    \end{notation}
\end{dang}
\begin{dang}
    \begin{definition}
        Let $A\in\text{M}_{n}(R)$. Then $A$ is said to be \textit{unimodular} if exists $B\in\text{M}_n(R)$ such that $AB=BA=I$.
    \end{definition}
\end{dang}
Important examples of unimodular matrices, are the \textit{elementary} matrices. These matrices may be used to do elementary operations with the row and columns of another matrix. For example, let $A=(A_{ij})\in \text{M}_n(R)$ and $\sigma\in S_n$, define $\sigma(A)=(A_{\sigma(i)j})$. Particularly, $(\sigma(I)A)_{ij}=\sum_{k=1}^n\delta_{\sigma(i)k}A_{kj}=A_{\sigma(i)j}$. So $\sigma(I)A=\sigma(A)$. This means that multiplication from the left on $A$ by $\sigma(I)$ has the effect of permuting the rows of $A$ by $\sigma$. Now, define $Y_{ij}(\alpha)=I+\alpha E_{ij}$ for $i\neq j$. Then, if $A\in\text{M}_n(R)$ is a matrix whose rows are $\{R_1,\dots,R_n\}$, then multiplication from the left on $A$ by $Y_{ij}(\alpha)$ has the effect of substituting the row $R_i$ by $R_i+\alpha R_j$. Clearly $\sigma(I)\sigma^{-1}(I)=\sigma^{-1}(I)\sigma(I)=I$ and $Y_{ij}(\alpha)Y_{ij}(-\alpha)=Y_{ij}(-\alpha)Y_{ij}(\alpha)Y=I$. So these matrices are unimodular.
\begin{dang}
    \textbf{Observation:} We can use elementary matrices $E$ to do elementary operations on columns of a matrix $A$. We just need to transpose $E$ and multiply on the right of $A$. See example~\ref{example_for_invariant_factor_theorem_1}.
\end{dang}
\begin{dang}
    \begin{definition}\cite{curtis1966representation}
        For $A,B\in\text{M}_{mn}(R)$, we say that $A$ and $B$ are \textit{equivalent} if there exists $X_1\in\text{M}_m(R)$ and $X_2\in\text{M}_n(R)$ such that $B=X_1AX_2$. In this case we will write $A\sim B$.
    \end{definition}
\end{dang}
\begin{dang}
    \begin{theorem}(Invariant factor theorem for matrices)\label{invariant_factors_theorem}\cite{curtis1966representation}
        Let $R$ be a principal ideal domain and $A$ a rectangular matrix with entries in $R$. Then $A\sim$diag$(\delta_1,\delta_2,\dots,\delta_k,0,\dots,0)$ where $\delta_i\in R-\{0\}$ and $\delta_i|\delta_{i+1}\;\forall i=1,\dots,k-1$. The elements $\delta_1,\dots,\delta_k$ are called \textit{invariant factors} and are uniquely determined up to unit factors.
    \end{theorem}
\end{dang}
Rather than reproduce the proof of this theorem, we offer examples that illustrate the general procedure for determining equivalence. 

\begin{dang}
    \begin{definition}\cite{hungerford2012algebra}
        A \textit{unique factorization domain} (UFD) is a ring $R$ such that every nonzero element can be written uniquely, up to unit factors, as a product of primes.
    \end{definition}
\end{dang}

\begin{dang}
    \begin{example}\label{example_for_invariant_factor_theorem_1}
        Consider
        \begin{equation}
            A = \begin{pmatrix}
                2 & 0 & 1\\
                0 & 2 & 1\\
                0 & 0 & -1
            \end{pmatrix}\in M_3(\bb{Z})
        \end{equation}
        This matrix is equivalent to diag$\{1,2,2\}$.
    \end{example}
\end{dang}
Multiply $A$ on the right by $(13)(I), Y_{21}(-1)$ and on the left $Y_{31}(1)$ to get
\begin{equation}
    \begin{pmatrix}
        1 & 0 & \;\;\;2\\
        0 & 2 & -2\\
        0 & 0 & \;\;\;2
    \end{pmatrix}
\end{equation}
Now, multiply this matrix on the right by $Y_{31}(-2)^T$ and $Y_{32}(1)^T$ to obtain diag$\{1,2,2\}$.\par
The next example illustrates the case where in some step, there is one element that doesn't divide some other element in the matrix. In this case the strategy will be the following: let $A=(a_{ij})$, where for some $(i,j)$ there is $(m,n)$ such that $a_{ij}$ does not divide $a_{mn}$. Using elementary operations on rows and columns we can assume, without loss of generality, that $(i,j)=(1,1)$ and $(m,n)=(1,2)$ or $(2,1)$. Suppose the former holds. Since $R$ is a principal ideal domain, we can compute the gcd$(a_{11},a_{12})=\beta$. In this case, there are $\gamma_1,\gamma_2\in R$ such that $\gamma_1a_{11}+\gamma_2a_{12}=\beta$. Also, since $\beta | a_{1i}, i=1,2$, there exist $\epsilon_i, i=1,2$ such that $\epsilon_1\gamma_1-\epsilon_2\gamma_2=1$. In this case, define
\begin{equation}
    X = \begin{pmatrix}
        \gamma_1 & \epsilon_2 & \\
        \gamma_2 & \epsilon_1 &  \\
                 &            & I
    \end{pmatrix}
\end{equation}
This matrix is unimodular with inverse given by
\begin{equation}
    X^{-1} = \begin{pmatrix}
        \;\;\epsilon_1 & -\epsilon_2 & \\
        -\gamma_2 & \;\;\epsilon_1 &  \\
                 &            & I
    \end{pmatrix},
\end{equation}
So $AX\sim A$ and $(AX)_{11}=\beta$. We are working in $R$ a PID, it's known that every PID is a UFD, so each element has a factorization into a finite product of prime elements. Since $a_{11} $ does not divide $a_{12}$, $\beta$ has strictly fewer prime factors than $a_{11}$. Applying this reduction multiple times, we get at some step an element $\gamma$ at $(1,1)-$position that divide all the others, so we can apply elementary operations with rows and columns as in the example~(\ref{example_for_invariant_factor_theorem_1}) to get
\begin{equation}
    A\sim\begin{pmatrix}
        \gamma &  \\
               & B
    \end{pmatrix}
\end{equation}
Now we repeat this procedure with the smaller matrix $B$.
\begin{dang}
    \begin{example}\label{example_for_invariant_factor_theorem_2}
        Consider 
        \begin{equation}
            A = \begin{pmatrix}
                t & -1 & \;\;0\\
                0 &\;\;t & -1\\
                0 & \;\;0 & \;\;t
            \end{pmatrix}\in \text{M}_3(k[[t]])
        \end{equation}
        This matrix is equivalent to diag$(1,1,t^3)$.
    \end{example}
\end{dang}
Since $t\ndiv -1$, we write $\gamma_1 t+\gamma_2(-1)=1$ where $\gamma_1=1$ and $\gamma_2=t-1$. So we also have $\gamma_1\epsilon_1-\gamma_2\epsilon_2=1$ for $\epsilon_1 = t$ and $\epsilon_2=1$. Therefore, we take
\begin{equation}
    X= \begin{pmatrix}
        1 & 1 & 0 \\
        t-1 & t & 0\\
        0 & 0 & 1
    \end{pmatrix}
\end{equation}
So 
\begin{equation}
   A\sim AX =\begin{pmatrix}
       1&0&0\\
       t^2-t&t^2&-1\\
       0&0&t
   \end{pmatrix}\sim
   \begin{pmatrix}
       1&0&0\\
       0&t^2&-1\\
       0&0&t
   \end{pmatrix}
\end{equation}
Now, repeating exactly the same procedure with the smaller matrix 
$\begin{pmatrix}
    t^2&-1\\
    0&\;\;t
\end{pmatrix}$ we get 
\begin{equation}
    \begin{pmatrix}
    1 & 2t^2\\
    t^3-t&-t^3
\end{pmatrix}\sim
\begin{pmatrix}
    1 & 0\\
    0&-t^3
\end{pmatrix}
\end{equation}
Finally we conclude that $A\sim$diag$(1,1,-t^3)$. Maybe note that it's equivalent to diag$(1,1,t^3)$.\par

Now let $M$ be a finitely generated free module over a principal ideal domain $R$ and $N\subset M$ a submodule. Then by theorem~(\ref{submodule_ofa_free_module_is_free}) $N$ has a basis. If $\{m_1,\dots,m_k\}$ is a basis of $M$ and $\{n_1,\dots,n_s\}$ is a basis of $N$, then we can write 
\begin{equation}
    (n_1,\dots,n_s)=(m_1,\dots,m_k)A
\end{equation}
where $A\in \text{M}_{ks}(R)$. So multiplying $A$ from the left by a unimodular matrix is equivalent to a change of basis of $M$ and multiplication on the right is equivalent to a change of basis of $N$. In this case, theorem~(\ref{invariant_factors_theorem}) says that we can choose a basis of $\{m_1,\dots,m_k\}$ of $M$ such that $\{\delta_1m_1,\dots,\delta_sm_s\}$ is a basis for $N$, where $\delta_i|\delta_{i+1}$ for $i=1,\dots,s-1$. For the uniqueness of $\{\delta_i\}$, see \cite{curtis1966representation}.
\begin{dang}
    \begin{example}\label{Klein_4group_example}
        Let $K=\{1,X,Y,XY\}$ be the Klein $4-$group, the unique non cyclic group of order $4$. It's an abelian group, so it's also a $\bb{Z}-$module. In this case
        \begin{equation}
            K\cong \frac{\bb{Z}}{(2)}\oplus\frac{\bb{Z}}{(2)}
        \end{equation}
    \end{example}
\end{dang}
 To see this, let $F$ be the free $\bb{Z}-$module generated by $K$, so we have a surjective homomorphism
        \begin{equation}
             f:F\rightarrow K: \alpha+\beta X+\gamma Y+\theta XY\mapsto X^\beta Y^\gamma (XY)^\theta
        \end{equation}
        We are using summation notation on $F$ to differ from the multiplication notation on $K$ and avoid confusion on operations defined on each module. But by theorem~(\ref{submodule_ofa_free_module_is_free}), Ker$f$ is a free submodule of $F$ and in fact, has a basis $\{1,2X,2Y,X+Y-XY\}$. We write
        \begin{equation}
            (1,2X,2Y,X+Y-XY)=(1,X,Y,XY)\begin{pmatrix}
                1 & 0 & 0 & \;\;0 \\
                0 & 2 & 0 & \;\;1\\
                0 & 0 & 2 & \;\;1\\
                0 & 0 & 0 & -1
            \end{pmatrix}
        \end{equation}
        By example~(\ref{example_for_invariant_factor_theorem_1}) we have
        \begin{equation}
            \begin{pmatrix}
                1 & 0 & 0 & \;\;0 \\
                0 & 2 & 0 & \;\;1\\
                0 & 0 & 2 & \;\;1\\
                0 & 0 & 0 & -1
            \end{pmatrix}\sim\text{diag}(1,1,2,2)
        \end{equation}
        Therefore, there is a basis of $\{m_1,m_2,m_3,m_4\}$ of $M$ such that $\{m_1,m_2,2m_3,2m_4\}$ is a basis of $N$. And we conclude that
        \begin{equation}
            K\cong\frac{F}{\text{Ker}f}\cong \frac{\bb{Z}m_1}{\bb{Z}m_1}\oplus\frac{\bb{Z}m_2}{\bb{Z}m_2}\oplus\frac{\bb{Z}m_3}{\bb{Z}2m_3}\oplus\frac{\bb{Z}m_4}{\bb{Z}2m_4}\cong\frac{\bb{Z}}{2\bb{Z}}\oplus\frac{\bb{Z}}{2\bb{Z}}
        \end{equation}
Of course $\{m_1,m_2,m_3,m_4\}$ can be found explicitly. By example~(\ref{example_for_invariant_factor_theorem_1}) we have 
\begin{equation}
    (m_1,m_2,m_3,m_4)=(1,X,Y,XY)Y_{31}(-1)Y_{21}(1)=(1,X+Y-XY,Y, XY)
\end{equation}
\begin{dang}
    \begin{example}\label{nilpotent_matrix_example}
        Let $\bb{F}$ be a field and 
        \begin{equation}
            A = \begin{pmatrix}
                0&1&0\\
                0&0&1\\
                0&0&0
            \end{pmatrix}\in\text{M}_3(\bb{F})
        \end{equation}
        Since $A$ is nilpotent, the vector space $M=\bb{F}^3$ becomes an $\bb{F}[[t]]-$module with action $t\cdot v\bydef Av$ for $v\in M$. In this case
        \begin{equation}
            M \cong\frac{\bb{F}[[t]]}{(t^3)}
        \end{equation}
    \end{example}
\end{dang}
To see this, let $F$ be the free $\bb{F}[[t]]-$module generated by the canonical basis $\{e_1,e_2,e_3\}$ of $F$. To avoid confusion, in $M$ we will use the letter $f_i, i=1,2,3$ for this basis. So, we have a surjective homomorphism of modules
        \begin{equation}
            f:F\rightarrow M: f_i\mapsto e_i
        \end{equation}
        In this case $\text{ker}(f)$ have basis $\{tf_1,tf_2-f_1,tf_3-f_2\}$ and
        \begin{equation}
            (tf_1,tf_2-f_1,tf_3-f_2)=(f_1,f_2,f_3)\begin{pmatrix}
                t&-1&0\\
                0&t&-1\\
                0&0&t
            \end{pmatrix}
        \end{equation}
        By example~(\ref{example_for_invariant_factor_theorem_2}), we have 
        \begin{equation}
            \begin{pmatrix}
                t&-1&0\\
                0&t&-1\\
                0&0&t
                 \end{pmatrix}\sim\text{diag}(1,1,t^3)
        \end{equation}
        Therefore, there is a basis $\{m_1,m_2,m_3\}$ of $F$ such that $\{m_1,m_2,t^3m_3\}$ is a basis of $\text{Ker}(f)$ and we conclude that 
        \begin{equation}
            M\cong \frac{F}{\text{Ker}(f)}\cong\frac{\bb{F}[[t]]}{(t^3)}
        \end{equation}
These examples illustrate the structure theorem for finitely generated modules over a principal integral domain.
\begin{dang}
\begin{theorem}\label{structure_theorem_for_modules_over_PID}
    \cite{curtis1966representation}
    Let $M$ be a finitely generated module over a principal ideal domain $R$. Then, there exist elements $\delta_i \in R - \{0\}$, uniquely defined up to unit factors, for $i = 1, \dots, k$, such that $\delta_i \mid \delta_{i+1}$, and
    \begin{equation}
        M \cong \frac{R}{R\delta_1} \oplus \dots \oplus \frac{R}{R\delta_k} \oplus R \oplus \dots \oplus R.
    \end{equation}
    The summand $R \oplus \dots \oplus R$ will be called the \textit{free} part of $M$.
\end{theorem}
\end{dang}
The general strategy for the proof closely follows the procedure demonstrated in the examples~(\ref{Klein_4group_example}) and~(\ref{nilpotent_matrix_example}). In the case that each $\delta_i$ factors as a products of primes, we can refine this result. To do this, we use the next lemma.
\begin{dang}
    \begin{lemma}\label{Za x Zb is isomorphic to Zab}\cite{curtis1966representation}
        Let $R$ be a principal ideal domain and $\delta=p_1^{k_1}\cdots p_n^{k_n}$ where $p_1,\dots,p_n$ are non-associated primes. Then, we have a module isomorphism
        \begin{equation}
            R/(\delta)\cong R/(p_1^{k_1})\oplus\cdots\oplus R/(p_n^{k_n})
        \end{equation}
    \end{lemma}
\end{dang}
Proof: Let $p,q$ be non-associated primes. Then $(p)\neq(q)\Rightarrow (p)+(q)=R$ because both $(p) \text{ and }(q)$ are maximal ideals (by theorem~(\ref{In_PID_every_primeIdeal_is_maximal})). Therefore, there exist $r,s\in R$ such that $rp+sq=1$. Now, define the map
\begin{equation}
    \varphi: R/(p)\oplus R/(q)\rightarrow R/(pq):(x+(p),y+(q))\mapsto py+qx+(pq)
\end{equation}
This map is clearly a well defined homomorphism of $R-$modules. Since $R/(pq)$ is cyclic and $\varphi(s+(p),r+(q))=1+(pq)$, we see that $\varphi$ is surjective. To see that it is injective let $\varphi(x+(p),y+(q))=0\iff py+qx=pqu$ for some $u\in R$. Now, multiplying $rp+sq=1$ by $y$, we get $yrp+ysq=y$ which implies that $y+rqx=pqu+ysq\Rightarrow y\in (q)$. Similarly, multiplying $rp+sq=1$ by $x$ we see that $x\in (p)$. 
This shows that $\varphi$ is injective and the rest of the proof follows by induction.

\begin{dang}
    \begin{example}
        For example, it is a consequence of the \textit{euclidean division algorithm} that the ring of integers $\bb{Z}$ is a unique factorization domain. The same goes for the ring of polynomials in one variable $\bb{F}[X]$ over a field $\bb{F}$.
    \end{example}
\end{dang}
It's possible to show, using the Zorn's lemma, that any principal ideal domain $R$ is a unique factorization domain (the converse doesn't holds, for example $\bb{Z}[X]$ is a unique factorization that is not a principal ideal) \cite{hungerford2012algebra}. But in this work, we will only work with the rings $\bb{Z}$ and $\bb{F}[X]$. These rings are both unique factorization domains (by the euclidean algorithm) and principal ideal domains. So, by this reason, we will state these two hypothesis together whenever we need them.
\begin{dang}
    \begin{corollary}\label{refinement_of_structure_theorem}\cite{curtis1966representation}
        If $R$ is a principal ideal domain. Then, any finitely generated $R-$module $M$ is of the form
        \begin{equation}
            M\cong \bigoplus_{p\text{ prime}}\bigoplus_{i\geq1}R/(p^{\mu_i(p)})
        \end{equation}
        where $\mu: p\in\{\text{primes elements of }M\}\mapsto \mu(p)=(\mu_1(p),\mu_2(p),\dots)\in\bb{N}^\infty$ is uniquely defined by $M$. Also, $\mu$ have the property that $\mu_j(p)=0$ if $j\geq j_p$ for some $j_p\geq1$ and $\mu(p)=0$ for all but a finite subset of primes elements $p$ of $R$.
    \end{corollary}
\end{dang}
This corollary will be useful in the classification of conjugacy classes of $GL_n(\fq)$ in the next chapter (see theorem~(\ref{cojugacy_classes_and_partition_functions})). 
\section{Localization}

\subsection{Localization of rings and its properties}
\begin{dang}
    \begin{definition}
    Let $R$ be a ring, $S\subset R$ is said to be a multiplicative system if $1\in S$ and whenever $x,y\in S\Rightarrow xy \in S$. So this is equivalent to say that $S$ is a multiplicative monoid.
    \end{definition}
\end{dang}

\begin{dang}
    \begin{definition}
        Let $R$ be a ring and $S\subset \text{Z}(R)$, the center of $R$, a multiplicative system. Define an equivalence relation on $R\times S$ by
        \begin{equation}
            (r_1,s_1)\sim   (r_2,s_2) \iff \exists\; s\in S\text{ such that } r_1s_2s=r_2s_1s
        \end{equation}
        We will write $r/s$ for the equivalence class containing $(r,s)$ under $\sim$.
    \end{definition}
\end{dang}

\begin{dang}
    \begin{theorem}\label{localization_produces_a_ring}\cite{rowen2008graduate}
    Let $R$ be a commutative ring and $S\subset R$ be a multiplicative system. Define 
    \[
    S^{-1}R\bydef (R\times S)/\sim
    \]
    where $\sim$ is as in the previous definition. Then $S^{-1}R$ is a ring with sum and multiplication defined by
    \begin{equation}
        \frac{x}{s}+\frac{y}{t}=\frac{xt+ys}{st}\text{ and }\frac{x}{s}\frac{y}{t}=\frac{xy}{st}
    \end{equation}
    \end{theorem}
\end{dang}
Proof: See~(\cite{rowen2008graduate}).
\begin{dang}
    \begin{lemma}
        The following map
        \begin{equation}
            f:R\rightarrow S^{-1}R:x\mapsto x/1
        \end{equation}
        is a ring homomorphism and it's called \textit{localization map}. In the case that $R$ is a principal ideal domain, $f$ is also injective.
    \end{lemma}
\end{dang}
\begin{dang}
    \begin{notation}
        For $A\subset R$ an arbitrary subset, we will write
        \begin{equation}
            S^{-1}A = \big\{x/s\in S^{-1}R\;\big|\;x\in A\text{ and }s\in S\big\}
        \end{equation}
    \end{notation}
\end{dang}
\begin{dang}
    \begin{lemma}\label{localization_of_preimages_of_ideals}
    Let $f: R\rightarrow S^{-1}R$ be as above. For any ideal $I\subset S^{-1}R$, then $f^{-1}(I)\subset R$ is an ideal. Moreover $I=S^{-1}f^{-1}(I)$.
    \end{lemma}
\end{dang}

Proof: To ease notation, write $\overline{I}\bydef f^{-1}(I)$. Let $x/s\in I$, then since $I$ is an ideal, we have $\frac{s}{1}\frac{x}{s}=\frac{x}{1}\in I\Rightarrow x\in\overline{I}$. So $x/s\in S^{-1}\overline{I}$ and $I\subset S^{-1}\overline{I}$. Conversely, if $x/s\in S^{-1}\overline{I}\Rightarrow x\in \overline{I}$ and $s\in S$. Moreover $x/1\in I\Rightarrow \frac{1}{s}\frac{x}{1}=\frac{x}{s}\in I$ and we get $S^{-1}\overline{I}\subset I$.
\begin{dang}
    \textbf{Remark: }Lemma ~(\ref{localization_of_preimages_of_ideals}) says that any ideal $I\subset S^{-1}R$ is of the form $S^{-1}\overline{I}$ for some ideal $\overline{I}\subset R$.
\end{dang}
 
Since $0\in S$ implies $S^{-1}R=(0)$, this case isn't interesting. Therefore we will be considering only multiplicative systems $S\subset R-\{0\}$.

\begin{dang}
    \begin{lemma}\label{R_is_PID_localization_PID}
    Let $R$ be an integral domain and $S\subset R-\{0\}$ a multiplicative system. Then $S^{-1}R$ is an integral domain.
    \end{lemma}
\end{dang}
Proof: Suppose that $\frac{r_1}{s_1}\frac{r_2}{s_2}=0$. Then, there exists $s\in S$ such that $r_1r_2s=0\Rightarrow r_1r_2=0$ because $s\neq0$. So $r_1=0$ or $r_2=0\Rightarrow r_1/s_1=0$ or $r_2/s_2=0$. 

\begin{dang}
\begin{definition}
    A ring $R$ is said to be \textit{local} if it has a unique maximal ideal.
\end{definition}
\end{dang}
\begin{dang}
    \begin{theorem}\label{localization_at_prime_ideals_on_a_PID}
      Let $R$ be a PID  and $P\subset R$ a prime ideal. Then $S\bydef R-P$ is a multiplicative system. We claim that $R_P\bydef S^{-1}R$ is a local ring.
    \end{theorem}
\end{dang}
Proof: The candidate for the maximal ideal is
\begin{equation}
    I \bydef S^{-1}P = \big\{x/s \in R_P \;\big|\;x\in P\text{ and }s\in S\big\} 
\end{equation}
First, let's show that $I$ is a maximal ideal. By lemma~(\ref{In_PID_every_primeIdeal_is_maximal}), we have that $P$ is a maximal ideal in $R$. Now, let $I\subset J$, $J$ being an ideal. Then, if $\bar{J}\bydef\{x\in R\;|\;x/1\in J\}\Rightarrow\bar{J}\subset R$ is an ideal. In particular  $P\subset\bar{J}\Rightarrow \bar{J}=P$. Then by lemma~(\ref{localization_of_preimages_of_ideals}) we have $J = I$ or $J = R_P$. Now, let $J$ be any proper ideal in $R_P$. Consider two cases: If $\bar{J}\cap S\neq\varnothing$, then there exist $x\in \bar{J}\cap S\Rightarrow x/1\in \bar{J}$. But $1/x\in R_P$ and since $J$ is an ideal, we have $\frac{1}{x}\frac{x}{1}=\frac{1}{1}\in J$ and $J=R_P$. So we exclude this case. Now, if $\overline{J}\cap S=\varnothing\Rightarrow \bar{J}\subset P$ and $J\subset I$. Therefore $I$ contains all the proper ideals of $R_P$.
\par
The definition of localization can be extended to modules.
\begin{dang}
    \begin{definition}
        Let $R$ be a commutative ring, $S\subset R$ a commutative system and $M$ an $R-$module. We define an equivalence relation on $M\times S$ by $(m,s)\sim(n,t)$ if $\exists \;u\in S$ such that $mtu=nsu$. We will write $m/s$ for the equivalence class containing $(m,s)$ in $M\times R$.
    \end{definition}
\end{dang}
\begin{dang}
    \textbf{Remark: }Since $R$ is considered to be a commutative ring, $M$ will be regarded as a $(R,R)-$bimodule. So we can write the $R$ action from either direction.
\end{dang}
\begin{dang}
    \begin{theorem}\cite{rowen2008graduate}
        Define 
        \begin{equation}
            S^{-1}M\bydef(S\times R)/\sim
        \end{equation}
        then $S^{-1}M$ is an $S^{-1}R-$module with sum defined by
        \begin{equation}
            \frac{m}{s}+\frac{n}{t} = \frac{mt+ns}{st}
        \end{equation}
        and $S^{-1}R-$action by
        \begin{equation}
            \frac{x}{t}\frac{m}{s}=\frac{xm}{ts}
        \end{equation}
    \end{theorem}
\end{dang}
Proof: The proof is completely analogous to that of theorem~(\ref{localization_produces_a_ring}).
\subsection{Functorial properties of localization}
Both the localization of rings and modules have functorial properties that we will be useful in the next section.
\begin{dang}
\begin{notation}\label{notation for x_M which maps m to xm}
    Let $M$ be an $R-$module. Each $x\in R$ defines a map $x_M:M\rightarrow M: m\mapsto xm$. So we have a homomorphism $(\cdot)_M:R\rightarrow \End{M}:x\mapsto x_M$. Also, for any subset $S\subset R$, we write $S_M$ for the set of endomorphisms $s_M\in\text{End}(M)$ such that $s\in S$. 
\end{notation}
\end{dang}
\begin{dang}
    \begin{theorem}\label{functorial_property_of_localization}
        Let $R$ be a commutative ring and $S\subset R$ a multiplicative system. Then the category of $S^{-1}R-$modules is equivalent to the category of $R-$modules $M$ with the property that $S_M\subset\text{Aut}(M)$.
    \end{theorem}
\end{dang}
Proof: Define the functor $F:\cat{R}\rightarrow\cat{S^{-1}R}$ that associates to each $R-$module $M$ the same set, but with $S^{-1}R-$action defined by $(x/s)\cdot m\bydef x_M\circ s_M^{-1}(m)$ for $m\in M, x\in R\text{ and }s\in S$. Now define another functor $G:\cat{S^{-1}R}\rightarrow\cat{R}$ that associates to each $S^{-1}R-$module $M$, the same set, but with action $x\cdot m\bydef (x/1)m$ for $m\in M\text{ and }x\in R$. 
\begin{dang}
    \begin{theorem}\label{functor_from_R_to_localization}
        Let $R$ be a commutative ring and $S\subset R$ a multiplicative system. Then, there is a functor 
        \begin{equation}
            G:\cat{R}\rightarrow\cat{S^{-1}R}: M\mapsto S^{-1}M
        \end{equation}
        that associates to each $f\in\text{Hom}_{R}(M,N)$ a $G(f)\in\text{Hom}_{S^{-1}R}(S^{-1}M,S^{-1}N)$, defined by $G(f)(x/s)\bydef f(x)/s$. In fact, is an exact functor.
    \end{theorem}
\end{dang}
Proof: The fact that it is a functor is clear, now lets show that this in fact, defines an exact functor. Let
\begin{center}
    \begin{tikzcd}
    L \arrow[r, "u"]& M \arrow[r, "v"]& N
    \end{tikzcd}
\end{center}
be an exact sequence. Since $F$ is a functor, we also have a sequence
\begin{center}
    \begin{tikzcd}
    S^{-1}L \arrow[r, "G(u)"]& S^{-1}M \arrow[r, "G(v)"]& S^{-1}N
    \end{tikzcd}
\end{center}
In fact, it is clear that we have $\Image{G(u)} \subset\Ker{G(v)}$ because $G(v)\circ G(u)=G(vu)=G(0)=0$. We claim that it is an exact sequence. Now, let $x/s\in\Ker{G(v)}$. Then $v(x)/s=0\Rightarrow(v(x),0)\sim (0,1)$. So, there is a $t\in S$ such that $v(x)t=0\Rightarrow v(xt)=0\Rightarrow xt\in\Ker{v}=\Image{u}$. So there exist $y\in L$ such that $u(y)=tx$ and $G(u)(y/t)=x/s$.
\begin{dang}
    \begin{corollary}\label{localization of a quotient if the quotient of localizations}
        Let $N\subset M$ be a submodule, then $S^{-1}M/S^{-1}N\cong S^{-1}(M/N)$. Particularly, if $I\subset R$ is an ideal, we have $S^{-1}(R/I)\cong S^{-1}R/S^{-1}I$. 
    \end{corollary}
\end{dang}
We saw two functors related to the localization map. If $S_M\subset\text{Aut}(M)$, the first functor takes an $R-$module $M$ to the same set, but viewed as an $S^{-1}R-$module. The second one takes an $R-$module $M$ to the $S^{-1}R-$module $S^{-1}M$ in general. If the first case applies, there is no difference between these two functors. This is the content of the next theorem.
\begin{dang}
    \begin{theorem}\label{localization_does_not_change_the_module}
        If $S_M\subset \text{Aut}(M)$, then $S^{-1}M\cong M$ as an $S^{-1}R-$modules (or equivalently, as $R-$modules by theorem~(\ref{functorial_property_of_localization})).
    \end{theorem}
\end{dang}
Proof: We claim that the localization map $f:M\rightarrow S^{-1}M:m\mapsto m/1$ is an isomorphism of $S^{-1}R-$modules. To show this, lets make an observation: For $m\in M$ and $s\in S$ we have $(s_M^{-1}(m))/1=m/s$. In fact $(s_M^{-1}(m),1)\sim(m,s)$. Therefore $(r/s)f(m)=(r/s)(m/1)=rm/s=(r_M\circ s_M^{-1}(m))/1=f((r/s)\cdot(m))$. Moreover, $f(s_M^{-1}(m))=m/s$ which implies that $f$ is surjective. To see injectivity, let $f(m)=0\Rightarrow (m,1)\sim(0,1)\iff$ 
there exists $s\in S$ such that $0=sm=s_M(m)$. But $S_M\subset \text{Aut}(M)$, which implies that $m=0$.  

\section{Modules over discrete valuation rings}
\subsection{Partitions of positive integers}
\begin{dang}
    \begin{definition}\cite{macdonald1998symmetric}
    A partition is any (finite or infinite) nonincreasing sequence $\lambda=(\lambda_1,\lambda_2,\dots)$ of positive integers such that $\lambda_i=0$ for $i\gg 1$. We choose not to distinguish between sequences that differ only by a string of zeros at the end. The numbers $\lambda_i$ are called the \textit{parts} of $\lambda$. 
\end{definition}
\end{dang}
We emphasize the fact we don't make any distinction between partitions like 
\[
(5,3,1),(5,3,1,0) \;\text{ or }\; (5,3,1,0,0,\dots).
\]


\begin{dang}
    \begin{definition}
    Following \cite{macdonald1998symmetric}, we will use the following characteristic of a partition of a positive integer:
        \begin{enumerate}
             \item$l(\lambda)$ is called the \textit{length} of $\lambda$, and is defined to be the number of nonzero parts in $\lambda$.
             \item $m_i(\lambda)$ is called \textit{multiplicity} of $i$ in $\lambda$, and is defined to be the number of parts of $\lambda$ that are equal to $i$.
             \item $|\lambda|$ is called the \textit{weight} of $\lambda$, and is defined to be the sum of its parts.
        \end{enumerate}
    \end{definition}
\end{dang}

\begin{dang}
    \begin{notation}
        Another common way to write partition is $\lambda=(1^{n_1}2^{n_2}\cdots)$. This notation emphasizes the multiplicities of each part of $\lambda$.
    \end{notation}
\end{dang}

\begin{dang}
    \begin{definition}
        Let $\lambda$ be a partition. If $|\lambda|=n$, we say that $\lambda$ is a partition of $n$. The set of all partitions of $n$ is denoted by $\mathcal{P}_n$ and the set of all partitions by $\mathcal{P}$. Its cardinality $p_n=|\mathcal{P}_n|$ can be computed by as the coefficient of $x^n$ in the following generating series
        \begin{equation}
            \sum_{i\geq1}p_ix^i = \prod_{j\geq1}\frac{1}{1-x^j}
        \end{equation}
    \end{definition}
\end{dang}
\begin{dang}
    \begin{definition}\cite{macdonald1998symmetric}
        Given a partition $\lambda$, we define its \textit{conjugate} $\lambda^\prime$, by defining $\lambda_i^\prime$ to be the number of parts of $\lambda$ that are greater than or equal to $i$. Since a partition is a nonincreasing sequence, it is clear that $\lambda^\prime$ is also a partition.
    \end{definition}
\end{dang}
\begin{dang}
    \textbf{Observation}: Another useful way to write the parts of $\lambda^\prime$ is by the formula $\lambda_i^\prime=\text{Max}\{k\geq1\;|\;\lambda_k\geq i\}$. In fact, let $a=\text{Max}\{k\geq1\;|\;\lambda_k\geq i\}$, then $\lambda_a\geq i$ and $\lambda_j\geq\lambda_a\;\forall j\;(1\leq j\leq a)$. But, since $\lambda_j<i$ for $j>a$, we have $\#\{k\geq1\;|\;\lambda_k\geq i\}=\#\{1,\dots,a\}=a$. 
\end{dang}
To each partition $\lambda$ we can associate a diagram as follows. The diagram of $\lambda$ is by definition the set $T_{\lambda}=\{(i,j)|\;j\leq\lambda_i\}$ where each pair $(i,j)\in T_\lambda$ is called a \textit{square}. Clearly, the $i$th row of $T_\lambda$ consist of $\lambda_i$ squares. To draw diagrams we use matrix convention, that is, rows going downward as $i$ increases, and columns indexed by $j$, going to the right as $j$ increases. With this convention, we claim that $T_{\lambda^\prime}=T_\lambda^t$, where $(-)^t$ denotes transpose, as for matrices. This follows from  $j\leq\lambda_i^\prime\iff \lambda_j\leq i$ because $\lambda_i^\prime = \text{Max}\{k\;|\;\lambda_k\geq i\}$. So $\lambda_i^\prime$ counts the number of squares in the $i$th column of the diagram $T_\lambda$.
\begin{dang}
    \begin{example}
        The diagram of the partition $\lambda=(2,1^2)$ is
        \begin{center}
            \ydiagram{2,1,1}
        \end{center}
    \end{example}
\end{dang}
\begin{dang}
    \begin{definition}\cite{macdonald1998symmetric}
        Given two partitions $\lambda$ and $\mu$, we write $\mu\subset \lambda$ if $\lambda_i-\mu_i\geq0$ for any $i\geq1$. So $\subset$ gives a  partial order in $\mathcal{P}$.
    \end{definition}
\end{dang}
\begin{dang}
    \begin{definition}\cite{macdonald1998symmetric}
        Let $\lambda,\mu\in\mathcal{P}$ such that $\mu\subset\lambda$. We define the skew diagram $\theta=\lambda-\mu$ to be the set theoretical difference between the diagrams of $\lambda$ and $\mu$. In this case, we will write $\theta_i=\lambda_i-\mu_i$ to be the number of squares on the $i$th lines on the skew diagram $\theta$. Also, we define the \textit{weight} of the skew diagram $\theta$ to be the sum $\sum_{i\geq1}\theta_i=|\lambda|-|\mu|$.
    \end{definition}
\end{dang}
We can see that if $\mu\subset\lambda$, then $\mu^\prime\subset\lambda^\prime$. Define the skew diagram $\theta=\lambda-\mu$. In this case we write $\theta^\prime=\lambda^\prime-\mu^\prime$. Another partial order on $\mathcal{P}$ that will be useful is the following.
\begin{dang}
    \begin{definition}
        Given $\mu,\lambda\in\mathcal{P}$, we define a partial order on $\mathcal{P}$ by $\mu\leq\lambda$ if 
        \[
        \mu_1+\cdots+\mu_i\leq \lambda_1+\cdots+\lambda_i\;\text{ for all}\; i\geq1.
        \]
    This is called the \textit{dominance} partial order in $\mathcal{P}$ \cite{macdonald1998symmetric}.
    \end{definition}
\end{dang}
\subsection{Discrete valuation rings}
\begin{dang}
    \begin{definition}
        A \textit{discrete valuation ring} $R$ is a principal ideal domain which is local with maximal ideal $P\neq(0)$. The \textit{residue field} of $R$ is defined by the quotient $R/P$.
    \end{definition}
\end{dang}
With this definition we can already see that there is just one simple $R-$module, namely the residue field $R/P$. In fact, if $M=Rm$ (any simple module must be cyclic), then we can define a surjective homomorphism $\varphi:R\rightarrow M:r\mapsto rm$. So $R/\text{Ker}(\varphi)\cong M$. But again, $M$ is simple, so $\text{Ker}(\varphi)$ must be a maximal ideal of $R$. Therefore $\text{Ker}(\varphi)=P$.\par
Next, an important property of discrete valuation rings.
\begin{dang}
    \begin{theorem}\label{Theorem: every element in a discrete valuation ring is upm}
        Let $R$ be a discrete valuation ring, and $p\in R$ a fixed generator of the unique maximal ideal. The, every $x\in R$ is of the form $x=up^m$ for some $m\geq0$ and unit $u\in R$. 
    \end{theorem}
\end{dang}
Proof: Since $R$ is noetherian, any non unit $a\in R$ is contained in $(p)$. Therefore, the set $R-(p)$ is precisely the set of units in $R$. We claim that for any $x\in R$ there exists a maximal $m\geq0$ such that $x\in (p^m)$. To show this, it's sufficient to prove that $\bigcap_{m\geq0}(p^m)=(0)$. By contradiction, suppose that there exists a nonzero element $x\in \bigcap_{m\geq0}(p^m)$. Then there is a sequence $(u_1,u_2,\dots)$ such that  $x=u_mp^m\;\forall m\geq1\Rightarrow u_m=pu_{m+1}\;\forall m\geq 1$. Therefore, we have a strictly increasing chain of ideals
\begin{equation}
    (u_1)\subset (u_2)\subset \cdots
\end{equation}
which contradicts the fact that $R$ is noetherian (see theorem~(\ref{principal_ideal_is_noetherian})).\par
Now, for $x\in R$, let $m\geq0$ be such that $x\in(p^m)$ and $x\notin(p^{m+1})$. Then, in this case $x=up^m$ where $u\notin(p)$, so $u$ is a unit.
\begin{dang}
    \begin{example}\label{residue_field_of_Rp}
       Let $R$ be a PID. By theorem~(\ref{localization_at_prime_ideals_on_a_PID}), $R_P$ is a discrete valuation ring. In this case, the residue field is given by $R_P/S^{-1}P\cong S^{-1}(R/P)$. Moreover if $S_R\subset\text{Aut}(R)$, then by the theorem~(\ref{localization_does_not_change_the_module}), the localization map $f:R\rightarrow S^{-1}R$ is an isomorphism of $S^{-1}R-$modules. But this map is also a ring homomorphism. Therefore, the residue field of $R_P$ is $S^{-1}(R/P)\cong R/P$.
    \end{example}
\end{dang}
\begin{dang}
    \begin{example}
        Another example of discrete valuation ring is given by the ring of Taylor series $\bb{F}[[t]]$. In this case, the unique maximal ideal is $(t)$ and the residue field is $\bb{F}[[t]]/(t)\cong \bb{F}$.
    \end{example}
\end{dang}
\subsection{Finite modules over a discrete valuation ring}
In this section, all rings $R$ considered are discrete valuation rings. Also, we will denote the maximal ideal by $P=(p)$ (so $p$ is a fixed generator). Now, if $M$ if a finitely generated $R-$module, then by theorem~(\ref{structure_theorem_for_modules_over_PID}) we have
\begin{equation}\label{decomposition_of_M_1}
    M\cong\left(\bigoplus_{i=1}^r\frac{R}{(p^{\lambda_i})}\right)\oplus F
\end{equation}
where $F$ is the free part of $M$ and $\lambda=(\lambda_1,\lambda_2,\dots)$ is a partition. Now we will be concerned with finitely generated modules $M$ over a discrete valuation ring $R$ that have the additional property that $p^rM=(0)$ for some $r\geq1$. In our context, these modules will be called \textit{finite} $R-$modules.
\begin{dang}
    \begin{definition}
        Let $M$ be a finitely generated module over a discrete valuation ring $R$ with maximal ideal $P=(p)$. Then $M$ is said to be a \textit{finite} $R-$module if $p^rM=(0)$ for some $r\geq1$. 
    \end{definition}
\end{dang}
In this case, the free part in the equation ~(\ref{decomposition_of_M_1}) must be $F=(0)$. So
\begin{equation}\label{decomposition_of_M_2}
     M\cong\bigoplus_{i=1}^r\frac{R}{(p^{\lambda_i})}
\end{equation}
Now, lets show that if $|R/(p)|=|\bb{F}|<\infty$, then these modules have a finite number of elements. 
\begin{dang}
    \begin{lemma}\label{quotient_of_finite_modules}
        Let $M$ be a module and suppose that there exists a finite submodule $N\subset M$ such that $M/N$ is of finite cardinality, then $M$ is of finite cardinality.
    \end{lemma}
\end{dang}
Proof: Let $N = \{n_1, \dots, n_l\}$, so that $[x] = \{x + n_1, \dots, x + n_l\}$. Now, if $M/N = \{[x_1], \dots, [x_k]\}$, then we have
\[
    M \subset \left\{ x_i + n_j \mid i = 1, \dots, k \text{ and } j = 1, \dots, l \right\}.
\]
\begin{dang}
\begin{theorem}
    If the residue field $\bb{F}$ of $R$ has $q$ elements, then $|R/(p^r)|=q^r$.
\end{theorem}
\end{dang}
Proof: By induction assume that the statement is true for $r-1$. Using the isomorphism theorem, we get
\begin{equation}
    \frac{R/(p^r)}{(p^{r-1})/(p^r)}\cong \frac{R}{(p^{r-1})}
\end{equation}
By lemma~(\ref{quotient_of_finite_modules}) $R/(p^r)$ is finite and
\begin{equation}
    \Bigg|\frac{R/(p^r)}{(p^{r-1})/(p^r)}\Bigg| = \frac{|R/(p^r)|}{|(p^{r-1})/(p^r)|}= \Bigg|\frac{R}{(p^{r-1})}\Bigg| = q^{r-1}
\end{equation}
But $(p^{r-1})/(p^r)\cong R/(p)$ and $|(p^{r-1})/(p^r)|=q$. So by the last equation $|R/(p^r)|=qq^{r-1}=q^r$.
\begin{dang}
    \begin{corollary}\label{cardinality_of_a_finite_Rmodule}
        Let $M$ be a finite $R-$module where the residue field of $R$ has $q$ elements. If $M$ has the decomposition as in the equation~(\ref{decomposition_of_M_2}), then
        \[
        |M|=q^{\sum_j\lambda_j}=q^{|\lambda|}
        \]
    \end{corollary}
\end{dang}
\subsection{Composition series}
\begin{dang}
    \begin{definition}\cite{erdmann2018algebras}
        Let $M$ be a module (over any ring). A \textit{series} for $M$ is a strictly decreasing sequence of submodules 
        \begin{equation}
            (0)=M_0\subset\dots\subset M_{n-1}\subset M_n = M
        \end{equation}
        If the additional property that each $M_{i}/M_{i-1}$ is simple, then we say that this series is a \textit{composition series} for $M$.
        In this case $n$ is called the \textit{length} of the composition series.
    \end{definition}
\end{dang}
Of course if $l(M)=0$, then $M=(0)$ and $l(M)=1$ if and only if $M\neq(0)$ is simple.
\begin{dang}
    \begin{definition}\cite{erdmann2018algebras}
        If $M$ has a composition series, we define $l(M)$ to be the least length of a composition series of $M$. If a module doesn't have a composition series, we put $l(M)=\infty$.
    \end{definition}
\end{dang}
\begin{dang}
    \begin{example}
        If $V$ is a finite dimensional vector space over $\bb{F}$, then $l(V)=\text{dim}_\bb{F}(V)$.
    \end{example}
\end{dang}
\begin{dang}
    \begin{example}    
    Consider the lattice of submodules of $\bb{Z}_{30}$.
    \begin{center}
        \begin{tikzcd}
                  & \bb{Z}_{30}=(1) \arrow[ld, dash]\arrow[d, dash]\arrow[rd, dash] & \\
                (2)\arrow[d, dash]\arrow[dr, dash] & (3)\arrow[dr, dash]\arrow[dl, dash] & (5) \arrow[dl, dash]\arrow[d, dash]\\
                (6)\arrow[dr, dash] & (10)\arrow[d, dash] & (15)\arrow[dl, dash]\\
                     & (0) & 
             \end{tikzcd}   
    \end{center}  
    Any path from the top to the bottom yields a composition series. For example
    \begin{equation}
        (0) \subset (6) \subset (3) \subset \bb{Z}_{30}
    \end{equation}
    In fact, since $(6) = \bb{Z}_{30}/6\bb{Z}_{30}\cong \bb{Z} /6\bb{Z}$ and $(3) = \bb{Z}_{30}/3\bb{Z}_{30}\cong \bb{Z}/3\bb{Z}$. We have $\bb{Z}_{30}/(3)\cong \bb{Z}_{10}$ and $(3)/(2)\cong\bb{Z}_{2}$ which are simple submodules of $\bb{Z}$.
    \end{example}
\end{dang}
\begin{dang}
    \begin{theorem}(Jordan-Holder theorem, \label{Jordan_Holder_theorem}\cite{erdmann2018algebras})
        Suppose that $M$ has a composition series of length $n$. Then, any other composition series of $M$ has length $n$. Also, any series of $M$ can be refined to yield a composition series.
    \end{theorem}
\end{dang}
Proof: See~(\cite{erdmann2018algebras}, Chapter III).
\par
An useful property of length, is that it is additive in the following sense.
\begin{dang}
    \begin{theorem}\cite{erdmann2018algebras}
        Suppose that we have a short exact sequence of modules of finite length
        \begin{center}
            \begin{tikzcd}
               0\arrow[r] & L \arrow[r, "u"]& M \arrow[r, "v"]& N \arrow[r]& 0
            \end{tikzcd}
        \end{center}
        then $l(M)=l(L)+l(N)$.
    \end{theorem}
\end{dang}
Proof: Let $l(L)=l$ and $l(N)=n$. By hypothesis, we have composition series
\begin{equation}
    \begin{split}
        (0) &= L_l \subset\dots\subset L_1 = L\\
        (0) &= N_n \subset\dots\subset N_1 = N
    \end{split}
\end{equation}
Since $u$ is injective, we also have another two series
\begin{equation}
\begin{split}
     (0) &= u(L_l) \subset\dots\subset u(L_1) = u(L)  \\
     v^{-1}(0) &= v^{-1}(N_n) \subset\dots\subset v^{-1}(N_1) = v^{-1}(N)
\end{split}    
\end{equation}
But $u(L)=\text{Ker}(v)=v^{-1}(0)$. And since $u$ is injective, the first series in the above equation is a composition series for $u(L)$. So we build a new series
\begin{equation}
     (0) = u(L_l) \subset\dots\subset u(L_1) = u(L)= v^{-1}(0) = v^{-1}(N_n) \subset\dots\subset v^{-1}(N_1) = v^{-1}(N) = M
\end{equation}
and we claim that this is a composition series for $M$. It's clear that if this is the case, then the length of this series is $n+l$. Therefore, from the above considerations, we just need to show that 
\begin{equation}
    \frac{v^{-1}(N_i)}{v^{-1}(N_{i-1})} 
\end{equation}
are simple modules. To do this, we will use the isomorphism theorem to show that 
\begin{equation}
     \frac{v^{-1}(N_i)}{v^{-1}(N_{i-1})} \cong \frac{v^{-1}(N_i)/\text{Ker}(v)}{v^{-1}(N_{i-1})/\text{Ker}(v)}\cong\frac{N_i}{N_{i-1}} 
\end{equation}
To do this, note that $N_i\cong v^{-1}(N_i)/\text{ker}(v)$ by the isomorphism theorem. So let $\varphi^\prime: N_i\rightarrow v^{-1}(N_i)/\text{Ker}(v):x\mapsto y + \text{Ker}(v)$ where $y\in v^{-1}(x)$ be an isomorphism.  To ease notation, lets write $A_i\bydef v^{-1}(N_i)/\text{Ker}(v)$. Now compose $\varphi^\prime$ with the projection $A_i\rightarrow A_i/A_{i-1}$ to get
\begin{equation}
    \varphi: N_i\rightarrow A_i/A_{i-1}: x\mapsto \overline{y}+A_{i-1}
\end{equation}
where $\overline{y}= y + \text{Ker}(v)$. Then $x\in\text{Ker}(\varphi)\iff\varphi(x)=\overline{y}+A_{i-1}=0\Rightarrow \overline{y}\in A_{i-1}$. This means that $y+\text{Ker}(v)\in v^{-1}(N_{i-1})/\text{Ker}(v)\iff y\in  v^{-1}(N_{i-1})\iff x\in N_{i-1}$. So, by the isomorphism theorem, we have $N_i/N_{i-1}\cong A_i/A_{i-1}$.
\begin{dang}
    \begin{corollary}
        Let $M$ be a module of finite length, if $N\subset M$ is of finite length, then $l(N)\leq l(M)$.
    \end{corollary}
\end{dang}
Now we will show that finite $R-$modules over discrete valuation rings have finite length.
\begin{dang}
  \begin{definition}\cite{curtis1966representation}
      Let $M$ be an $R-$module and $u\in M$. We write $\text{Ann}_R(u)$ (called the \textit{annihilator} of $u$) for the set of elements $r\in R$ such that $ru=0$. It's clear from the definition that this set forms an ideal in $R$.
  \end{definition}
\end{dang}
\begin{dang}
    \begin{lemma}\label{existence_of_primeideals_that_annihilate_u}
    Let $M\neq (0)$ be a module over a noetherian ring $R$. 
    Then there exists a prime ideal $P\subset R$ of the form $P=\text{Ann}_R(u)$ for some $u\in M$.
    \end{lemma}
\end{dang}
Proof: Let $\mathcal{P}=\{\text{Ann}_R(u)\;|\;u\in M-\{0\}\}$. Since $R$ is noetherian, $\mathcal{P}$ has a maximal element $P=\text{Ann}_R(u)$, i.e., does not exist $Q\in\mathcal{P}$ such that $Q\supset P,Q\neq P$. Note that $u\neq0\Rightarrow P\neq M$ (for example $1\notin P$). Now, $b\notin P\Rightarrow bu\neq0$ and $\text{Ann}_R(bu)\in\mathcal{P}$. But $\text{Ann}_R(u)\subset \text{Ann}_R(bu)\Rightarrow \text{Ann}_R(u)=\text{Ann}_R(bu)=P$ by maximality. So, if $ab\in P\Rightarrow abu=a(bu)=0$. Therefore $a\in \text{Ann}_R(bu)=P$. This shows that $P$ is prime. 
\begin{dang}
    \begin{lemma}\label{existence_of_submodule_such_that_the_quotient_is_simple}
    Let $M\neq0$, then there exists $M_1\subset M$ such that $M_1\cong R/P_1$ for some prime ideal $P_1$.
    \end{lemma}
\end{dang}
Proof: By lemma~(\ref{existence_of_primeideals_that_annihilate_u}), there exists a prime ideal of the form $P_1=\text{Ann}_R(u)$ for some $u\neq0$. Now, define the homomorphism $\varphi:R\rightarrow M: r\mapsto ru$. This map induces an injection $i:R/\text{Ker}(\varphi)\hookrightarrow M$. But $\text{Ker}(\varphi)=\text{Ann}_R(u)=P_1$. Therefore we have an inclusion $M_1\hookrightarrow M$ where $M_1\bydef i(R/P_1)\cong R/P_1$.
\begin{dang}
\begin{theorem}
    Let $M\neq(0)$ be a noetherian module over a noetherian ring $R$. Then $M$ has a series of the form
    \begin{equation}
         (0)\subset M_1\subset M_2\subset \dots\subset M_n= M
    \end{equation}
    such that $M_i/M_{i-1}\cong R/P_i$ for prime ideals $P_i\subset R$.
\end{theorem}
\end{dang}
Proof: By the lemma~(\ref{existence_of_submodule_such_that_the_quotient_is_simple}), there exists $M_1\subset M$ such that $M_1\cong R/P_1$ for some prime ideal $P_1\subset R$. If $M/M_1\neq (0)$, we can apply lemma~(\ref{existence_of_submodule_such_that_the_quotient_is_simple}) again to the quotient module $M/M_1$ to get a module $M_2^\prime\subset M/M_1$ such that $M_2^\prime\subset R/P_2$ for some prime ideal $P_2\subset R$. But $M_2^\prime=M_2/M_1$ for some $M_2$ containing $M_1$. Repeating this procedure, we get a chain of submodules: $(0)\subset M_1\subset M_2\subset \dots$ such that $M_i/M_{i-1}\cong R/P_i$. Now, since $M$ is noetherian, the chain must terminate at some $M_n= M$ (if $M_n\neq M$, then $M/M_n\neq(0)$ and we can apply the lemma~(\ref{existence_of_primeideals_that_annihilate_u}) again to get a new module $M_{n+1}$ strictly bigger than $M_n$). So we get
\begin{equation}
    (0)\subset M_1\subset M_2\subset \dots\subset M_n= M
\end{equation}
such that $M_i/M_{i-1}\cong R/P_i$ for prime ideals $P_i\subset R$.
\par
Now, if we add the condition that $R$ is a principal ideal domain, then each prime ideal $P_i$ is also maximal. Therefore $R/P_i$ are simple modules. So we get the following corollary.
\begin{dang}
    \begin{corollary}\label{a_finite_module_over_aDVR_have_finite_lenght}
        With the additional condition that $R$ is a principal ideal domain we have that $l(M)<\infty$.
    \end{corollary}
\end{dang}
Now, let $M$ be a finite $R-$module, where $R$ is a discrete valuation ring with finite residue field. Then, $R$ is noetherian because it's a principal ideal ring and $M$ is noetherian because it's of finite cardinality. 
\begin{dang}
    \begin{corollary}\label{finite_o_modules_have_finite_composition_series}
        If $M$ is a finite module over a discrete valuation ring, then $l(M)<\infty$.
    \end{corollary}
\end{dang}

\subsection{Formula for the number of automorphisms}
From now on, we will only be interested in the case that $R$ is a discrete valuation ring with finite residue field $\bb{F}=R/(p)$. The number of elements of $\bb{F}$ will be denoted by $q$.\par
In the equation~(\ref{decomposition_of_M_2}) it is not clear whether $\lambda$ is uniquely defined by $M$. We will see that this is the case. In fact, let $M$ be a module over a discrete valuation ring $R$ with the property that $p^{r}M=(0)$ and $p^{r-1}M\neq(0)$ for some $r\geq1$. We have a series
\begin{equation}\label{chain_of_modules_pM}
    (0)=p^{r}M\subset \dots \subset pM\subset M
\end{equation}
Since $p\left(p^{i-1}M/p^{i}M=(0)\right)$, each quotient $p^{i-1}M/p^{i}M$ can be viewed as a vector space over the residue field $\bb{F}=R/(p)$. 
\begin{dang}
    \begin{theorem}\label{dim_of_M_over_pM}\cite{macdonald1998symmetric}
        Let
        \begin{equation}
             M\cong\bigoplus_{i=1}^r\frac{R}{(p^{\lambda_i})}
        \end{equation}
        then we have $\text{dim}_{\bb{F}}(p^{i-1}M/p^{i}M)=\lambda_i^\prime$.
    \end{theorem}
\end{dang}
Proof: We have
\begin{equation}
\begin{split}
     p^{i-1}M &\cong \bigoplus_{\lambda_j\geq i}\frac{p^{i-1}R}{(p^{\lambda_j})} \\
              &= \bigoplus_{\lambda_j=i}\frac{p^{i-1}R}{(p^{\lambda_j})}\oplus\bigoplus_{\lambda_j\geq i+1}\frac{p^{i-1}R}{(p^{\lambda_j})}\\
              &= \bigoplus_{\lambda_j=i}\bb{F}\oplus\bigoplus_{\lambda_j\geq i+1}\frac{p^{i-1}R}{(p^{\lambda_j})}
\end{split} 
\end{equation}
So
\begin{equation}
\begin{split}
     p^{i-1}M/p^iM&\cong \bigoplus_{\lambda_j=i}\bb{F}\oplus\bigoplus_{\lambda_j\geq i+1}\frac{p^{i-1}R/(p^{\lambda_j})}{p^{i}R/(p^{\lambda_j})}\\
     &\cong \bigoplus_{\lambda_j=i}\bb{F}\oplus\bigoplus_{\lambda_j\geq i+1}\bb{F}\\
     &=\bigoplus_{\lambda_j\geq i}\bb{F}\\
     &=\bb{F}^{\lambda_i^\prime}
\end{split} 
\end{equation}
therefore $\text{dim}_{\bb{F}}(p^{i-1}M/p^{i}M)=\lambda_i^\prime$.\par
Since $\lambda$ is uniquely defined by the module $M$ we can formulate the following definition.
\begin{dang}
    \begin{definition}
        If $M$ has a decomposition as in equation~(\ref{decomposition_of_M_2}), then $M$ is said to have \textit{type} $\lambda$.
    \end{definition}
\end{dang}
Of course, two modules have the same type if and only if they are isomorphic. By corollary~(\ref{finite_o_modules_have_finite_composition_series}) we see that finite $R-$modules have finite length. Now we claim that if $M$ is a finite $R-$module of type $\lambda$, then $l(M)=|\lambda|$. In fact, take a composition series
\begin{equation}
    (0)=M_0\subset M_1\subset \dots \subset M_{l(M)}=M
\end{equation}
then each $M_i/M_{i-1}\cong\bb{F}$. In particular, $M$ is finite and
\begin{equation}
    |M| =  \prod_{i=1}^{l(M)}\Bigg|\frac{M_i}{M_{i-1}}\Bigg| = \prod_{i=1}^{l(M)}q=q^{l(M)}
\end{equation}
By corollary~(\ref{cardinality_of_a_finite_Rmodule}) we have that $q^{l(M)}=q^{|\lambda|}\Rightarrow l(M)=|\lambda|$. 
\begin{dang}
    \begin{notation}\cite{macdonald1998symmetric}
        For a partition $\lambda$, we denote $n(\lambda)=\sum_{i\geq1}(i-1)\lambda_i$. Summing the columns of $T_\lambda$ (with these numbers attached) we also see that 
        \[
        n(\lambda)=\sum_{i\geq1}\binom{\lambda_i^\prime}{2}
        \]
    \end{notation}
\end{dang}
Now let's prove the main theorem of this section.
\begin{dang}
    \begin{theorem}\label{numbero_of_automorphisms_of_M_of_type_lambda}\cite{macdonald1998symmetric}
        The number of automorphisms of a finite $R-$module of type $\lambda$ is
        \begin{equation}
            a_\lambda(q)=q^{|\lambda|+2n(\lambda)}\prod_{i\geq1}\varphi_{m_i(\lambda)}(q^{-1})
        \end{equation}
        where $\varphi_m(t)=(1-t)(1-t^2)\dots(1-t^m)$.
    \end{theorem}
\end{dang}
\begin{dang}
    \begin{definition}\cite{macdonald1998symmetric}
        Let $x\in M-\{0\}$, we shall say that $x$ has \textit{height} $r$ if $p^rx=0$ while $p^{r-1}x\neq0$. So $x$ has height $r$ if and only if $Rx\cong R/(p^r)$.
    \end{definition}
\end{dang}

\begin{dang}
    \begin{notation}
        Let $M$ be a finite $R-$module. We will denote $M_r$ the submodule of $M$ generated by elements of height less than or equal to $r$.
    \end{notation}
\end{dang}

To count the number of automorphisms of a finite $R-$module of type $\lambda$, we need to list the number of 
sequences $(x_1,x_2,\dots)$ such that each $x_i$ has height $\lambda_i$. So we proceed as follows: to count the number of 
elements $x_1\in M$ with height $\lambda_1$, we need to calculate the cardinality of the set $M_{\lambda_1}-M_{\lambda_1-1}$. 
Now, choose $x_1$ in this set, then we need to count the number of elements $x_2\in M$ of height $\lambda_2$ that don't belong to the 
submodule $Rx_1$. Repeating this procedure, we see that at each step we need to choose a $x_j$ such that 
$Rx_j\cap (Rx_{j-1}\oplus\dots\oplus Rx_1)=(0)$ where each $x_i$ has height $\lambda_i$. To count this, 
the only property of the submodule $N=Rx_{j-1}\oplus\dots\oplus Rx_1$ that we need, is that it is a submodule 
generated by elements of height greater than or equal to $\lambda_{j-1}$.

\begin{dang}
    \begin{lemma}\label{value_of_l(M_r)}\cite{macdonald1998symmetric}
        Let $M$ be a $R-$module of type $\lambda$. Then $l(M_r)=\lambda_1^\prime+\dots+\lambda_r^\prime$.
    \end{lemma}
\end{dang}
Proof: By theorem~(\ref{dim_of_M_over_pM}), we have
\begin{equation}\label{telescopic_sum_for_the_chain_pM}
\begin{split}
     \lambda_1^\prime+\dots+\lambda_r^\prime &=\sum_{i=1}^r\text{dim}_\bb{F}\left(\frac{p^{i-1}M}{p^iM}\right)\\
     &=\sum_{i=1}^r l\left(\frac{p^{i-1}M}{p^iM}\right)\\
     &=\sum_{i=1}^r (l(p^{i-1}M)-l(p^{i}M))\\
     &=l(M)-l(p^rM)
\end{split}
\end{equation}
But, by definition $M_r$ is the kernel of the map $M\rightarrow p^rM:m\mapsto p^rm$. Therefore $M/M_r\cong p^rM$ and by the additive property of length we get $l(M_r)=l(M)-l(p^rM)= \lambda_1^\prime+\dots+\lambda_r^\prime$.
\begin{dang}
    \begin{lemma}\cite{macdonald1998symmetric}
    Let $M$ be a finite $R-$module and $N\subset M$ a submodule generated by elements of height greater than or equal to $r$. Then $N_1\subset p^{r-1}N$.
    \end{lemma}
\end{dang}
Proof: Since $M$ is a finite $R-$module, $N$ is a finite $R-$module and $N=\bigoplus_i Rx_i$ where each $x_i\in M$ has height greater than or equal to $r$. Now, if $z\in N_1,z\neq0$ we have $z=\sum_ir_ix_i$. Moreover $pz=0\Rightarrow pr_ix_i=0$ for each $i$. Therefore $pr_i\in(p^r)\Rightarrow r_i=p^{r-1}s_i$ for each $i$. This shows that $z=p^{r-1}\sum_is_ix_i\in p^{r-1}N$.
\begin{dang}
    \begin{lemma}\label{x_has_height_r_and_RxcapN=0_iff}\cite{macdonald1998symmetric}
        Let $M$ be a finite $R-$module and $N\subset M$ a submodule generated by elements of height greater than or equal to $r$, and let $x\in M$. Then 
        \begin{equation}
           x\text{ has height }r\text{ and }Rx\cap N=(0) \iff x \in M_r-(M_{r-1}+N_r)
        \end{equation}
    \end{lemma}
\end{dang}
Proof: Suppose that $x$ has height $r$ and $Rx\cap N=(0)$. Since $x$ has height $r$, we have $x\in M_r$. Also, if $x\in M_{r-1}+N_r\Rightarrow p^{r-1}x\in N-\{0\}$ contradicting the hypothesis that $Rx\cap N=(0)$. Now if $x\in M_r-(M_{r-1}+N_r)$, we have that $x$ has height $r$. But, if $Rx\cap N\neq(0)\Rightarrow$ there exists $m<r$ such that $p^mx\in N-\{0\}\text{ (see theorem~(\ref{Theorem: every element in a discrete valuation ring is upm})})\Rightarrow p^{r-1}x\in N_1\subset N$. By the previous lemma, $p^{r-1}x=p^{r-1}y$ for some $y\in N$ and since $p^{r-1}x\neq0$ we have that $y\in N_r$. Therefore $p^{r-1}(x-y)=0\Rightarrow x-y\in M_{r-1}$ while $y\in N_r\Rightarrow\in x M_{r-1}+N_r$.
\begin{dang}
    \begin{lemma}\label{use_this_corollary_to_count_automorphisms}\cite{macdonald1998symmetric}
       Let $M$ be a finite $R-$module of type $\lambda$. If $N\subset M$ is a submodule of type $\nu$, then the number of elements $x\in M$ with height $r$ and satisfying $Rx\cap N=(0)$ is
        \begin{equation}
            q^{\lambda_1^\prime+\dots+\lambda_r^\prime}(1-q^{\nu_r^\prime-\lambda_r^\prime})
        \end{equation}
    \end{lemma}
\end{dang}
Proof: By lemma~(\ref{x_has_height_r_and_RxcapN=0_iff}) we just need to compute $|M_r-(M_{r-1}+N_r)|$. But
\begin{equation}
        |M_r-(M_{r-1}+N_r)|  = |M_r|-|M_{r-1}+N_r| = q^{l(M_r)}-q^{l(M_{r-1}+N_r)}
\end{equation}
and 
\begin{equation}
    \begin{split}
        l(M_{r-1}+N_r) &=l(M_{r-1})-l\left(\frac{M_{r-1}+N_r}{M_{r-1}}\right)\\
        &=l(M_{r-1})-l\left(\frac{N_r}{M_{r-1}\cap N_r}\right)\\
        &=l(M_{r-1})-(l(N_r)-l(N_{r-1}))\\
        &=\lambda_1^\prime+\dots\lambda_{r-1}^\prime - \nu_r^\prime\;\text{ (by lemma~(\ref{value_of_l(M_r)})})
    \end{split}
\end{equation}
Using the last two equations, we get $  |M_r-(M_{r-1}+N_r)|=q^{\lambda_1^\prime+\cdots+\lambda_{r}^\prime}(1-q^{\lambda_r^\prime-\nu_r^\prime})$
\par
Going back to the theorem~(\ref{numbero_of_automorphisms_of_M_of_type_lambda}), we see that each submodule $Rx_{j-1}\oplus\dots\oplus Rx_1$ has type $\nu=(\lambda_1,\dots,\lambda_{j-1})$. Therefore, using lemma~(\ref{use_this_corollary_to_count_automorphisms}) we see that
\begin{equation}
    \begin{split}
        |\text{Aut}(M)|&=\prod_{i\geq1}q^{\lambda_1^\prime+\dots+\lambda_{\lambda_i}^\prime}(1-q^{\nu_{\lambda_i}^\prime-\lambda_{\lambda_i}^\prime})\\
        &=\prod_{i\geq1}q^{\lambda_1^\prime+\dots+\lambda_{\lambda_i}^\prime}(1-q^{i-1-\lambda_{\lambda_i}^\prime})\; \text{ (because }\nu_{\lambda_i}^\prime=i-1)
    \end{split}
\end{equation}
Now, for $a\in \Image{\lambda}$ (the set of parts of $\lambda$), define $j_a\bydef\text{Min}\{j\geq1\;|\;\lambda_j=a\}$. Then $j_a=1+\lambda_a^\prime-m_a(\lambda)$. To verify this, just note that $m_a(\lambda)=\lambda_a^\prime-\lambda_{a+1}^\prime$ and $\lambda_{a+1}^\prime=\text{Max}\{j\geq1\;|\;\lambda_j\geq a+1\}=j_a-1$. Using this notation, we can write
\begin{equation}
    \begin{split}
        \prod_{i\geq1}q^{\lambda_1^\prime+\dots+\lambda_{\lambda_i}^\prime}(1-q^{i-1-\lambda_{\lambda_i}^\prime}) &=\prod_{a\in\Image{\lambda}}\prod_{i=0}^{m_a(\lambda)-1}q^{\lambda_1^\prime+\dots+\lambda_a^\prime}(1-q^{(i+j_a)-1-\lambda_a^\prime})\\
        &=\prod_{a\in\Image{\lambda}}\prod_{i=0}^{m_a(\lambda)-1}q^{\lambda_1^\prime+\dots+\lambda_a^\prime}(1-q^{i-m_a(\lambda)})\\
        &=\prod_{a\in\Image{\lambda}} q^{m_a(\lambda)(\lambda_1^\prime+\dots+\lambda_a^\prime)}\varphi_{m_a(\lambda)}(q^{-1})\\
        &=q^{\sum_{a\geq1}m_a(\lambda)(\lambda_1^\prime+\dots+\lambda_a^\prime)}\prod_{a\geq1}\varphi_{m_a(\lambda)}(q^{-1})\\
         &=q^{\sum_{a\geq1}(\lambda_a^\prime-\lambda_{a+1}^\prime)(\lambda_1^\prime+\dots+\lambda_a^\prime)}\prod_{a\geq1}\varphi_{m_a(\lambda)}(q^{-1})\\
         &=q^{\lambda_1^{\prime\;2}+\dots+\lambda_{\lambda_1}^{\prime\;2}}\prod_{a\geq1}\varphi_{m_a(\lambda)}(q^{-1})\\
            &=q^{|\lambda|+2n(\lambda)}\prod_{a\geq1}\varphi_{m_a(\lambda)}(q^{-1})
    \end{split}
\end{equation}
\section{Hall Algebras}
\subsection{Definition of the Hall algebra}
\begin{dang}
    \begin{definition}\cite{macdonald1998symmetric}
        Let $M$ be a finite $R-$module. Then $N\subset M$ is said to have \textit{cotype} $\mu$ in $M$ if $M/N$ is of type $\mu$.
    \end{definition}
\end{dang}
\begin{dang}
\begin{definition}\cite{macdonald1998symmetric}
     A finite $R-$module $M$ is called \textit{elementary} if $pM=(0)$. In this case we see that $M$ is elementary if and only if it is of type $(1^m)$ for some $m\geq1$.
\end{definition}
\end{dang}
An important property of elementary submodules is that they are precisely vector spaces over $\bb{F}$. In fact, if $pM=(0)$, the $R-$action pass through the quotient $(x+(p))\cdot m \bydef xm$ and $M$ becomes a vector space over $\bb{F}$. Conversely, if $M$ is a vector space over $F$, then we can compose the $\bb{F}-$action with the projection $R\rightarrow R/(p)=\bb{F}$ to get a $R-$module. 
\begin{dang}
    \begin{definition}\cite{macdonald1998symmetric}
        The \textit{Hall algebra} $H(R)$ is defined to be the free $\bb{Z}-$module on a basis $(u_\lambda)$ generated by all partitions $\lambda$. The product is defined by the equation 
        \begin{equation}
            u_\mu u_\nu\bydef\sum_\lambda G_{\mu\nu}^\lambda u_\lambda
        \end{equation}
        where the structure constants $G_{\mu\nu}^\lambda$ are defined as follows: Fix a finite $R-$module $M$ of type $\lambda$. The number $G_{\mu\nu}^\lambda$ is the number of chains
        \begin{equation}
            (0)\subset N\subset M
        \end{equation}
        such that $N$ is of type $\nu$ and cotype $\mu$ in $M$.
    \end{definition}
\end{dang}
\subsection{Calculantion of the structure constants}
From the additive property of the length, we get that $G_{\mu\nu}^\lambda\neq0$ only if $|\lambda|=|\mu|+|\nu|$. Before we prove the main properties of the Hall algebra, lets compute its most basics structure constants.
\begin{dang}
    \begin{lemma}\label{number_of_linearly_independent_vectors_inV}
        Let $V$ be a vector space over a finite field $\bb{F}$ of cardinality $q$. Then, for a fixed $i\leq d=\text{dim}_{\bb{F}}(V)$, the number of  sequences $(v_1,\dots,v_i)$ of linearly independent vectors in $V$ is
        \begin{equation}
             (q^d-1)(q^d-q)\dots(q^d-q^{i-1})
        \end{equation}
    \end{lemma}
\end{dang}
Proof: If $i=1$ we just need to choose a nonzero vector $v_1$. The number of such vectors is $q^d-1$, so the result is true for $i=1$. Now, assume by induction that the result is true for some $i\geq1$. Then, there are $(q^d-1)(q^d-q)\dots(q^d-q^{i-1})$ possibilities to choose a linearly independent sequence $(v_1,\dots,v_i)$ in $V$. Since the map $(a_1,\dots,a_i)\in\fq^i\mapsto \sum_{j=1}^ia_jv_j\in V$ is injective, the span of the sequence $(v_1,\dots,v_i)$ have exactly $q^i$ elements. So, there are $q^d-q^i$ ways to choose $v_{i+1}$ such that $(v_1,\dots,v_i,v_{i+1})$ is linearly independent. This shows that the number of linearly independent sequences of $i+1$ vectors is $(q^d-1)(q^d-q)\dots(q^d-q^{i-1})(q^d-q^i)$ which completes the induction argument.
\begin{dang}
    \begin{lemma}\label{number o LI sequence modulo N}
        Let $N\subset M$ be such that $pM\subset N$. Then, $M/N$ is a vector space of dimension $d=\text{dim}_{\bb{F}}(M/N)$, where $\bb{F}$ is the residue field of $R$. In this case, the number of sequences $(x_1,\dots,x_i)$ in $M$, for $i\leq d$, which are linearly independent modulo in $M/N$ is
        \begin{equation}
            (q^{d+n}-q^n)(q^{d+n}-q^{n+1})\dots(q^{d+n}-q^{n+i-1})
        \end{equation}
        where $n=l(N)$.
    \end{lemma}
\end{dang}
Proof: Fix $i\leq1$ A sequence $(x_1,\dots,x_i)$ in $M$ is linearly independent modulo $N$ if and only if $(\overline{x_1},\dots,\overline{x_i})$ is linearly independent in $M/N$.  By the previous lemma, there are $(q^d-1)(q^d-q)\dots(q^d-q^{i-1})$ linearly independent sequences $(\overline{x_1},\dots,\overline{x_i})$ in $V=M/N$. Now, $(\overline{x_1},\dots,\overline{x_i})=(\overline{y_1},\dots,\overline{y_i})$ if and only if $x_j-y_j\in N$. Since $N$ has $q^n$ elements, the number of sequences $(x_1,\dots,x_i)$ that have the same image under the projection is $q^{ni}$. Therefore, the number of sequences $(x_1,\dots,x_i)$ in $M$ that are linearly independent modulo $N$ is 
\begin{equation}
    q^{ni}(q^d-1)(q^d-q)\dots(q^d-q^{i-1})=(q^{d+n}-q^n)(q^{d+n}-q^{n+1})\dots(q^{d+n}-q^{n+i-1}).
\end{equation}
Lets compute explicitly a few values of $G_{\mu\nu}^\lambda$. It's convenient to introduce the Gauss polynomials $\smallgausspol{a}{b}_{t}=\prod_{i=1}^b(1-t^{a+i-1})/(1-t^i)$.
\begin{dang}
    \begin{example}
       $G_{(1^r)(1^s)}^{(1^{r+s})}$ counts the number of submodules $N\subset \bb{F}^{r+s}$ of type $(1^s)$ and cotype $(1^r)$. We claim that 
       \begin{equation}
           G_{(1^r)(1^s)}^{(1^{r+s})}=q^{sr}\gausspol{r+s}{s}_{q^{-1}}
       \end{equation}
    \end{example}
\end{dang}
In fact, these modules are precisely the subspaces of dimension $s$ in $\bb{F}^{r+s}$. Let $d=r+s$, from lemma~(\ref{number_of_linearly_independent_vectors_inV}) 
we get that there exists 
       \begin{equation}
           (q^d-1)(q^d-q)\cdots(q^d-q^{s-1})
       \end{equation}
       linearly independent sequences in $(v_1,\dots,v_s)$ in $M$. Each one of these sequences generates a subspace $N$ of dimension $s$. But, given $N$, it has precisely 
       \begin{equation}
           (q^s-1)(q^s-q)\cdots(q^s-q^{s-1})
       \end{equation}
       basis, i.e. linearly independent sequences which generated the whole $N$. Therefore, the number of subspaces of dimension $s$ in $M$ is the quotient
       \begin{equation}
       \begin{split}
              &\frac{(q^d-1)(q^d-q)\cdots(q^d-q^{s-1})}{(q^s-1)(q^s-q)\cdots(q^s-q^{s-1})}\\
              &=\frac{q^{ds}}{q^{ss}}\frac{(1-q^{-d})(1-q^{-d+1})\cdots(1-q^{-d+s-1})}{(1-q^{-s})(1-q^{-s+1})\cdots(1-q^{-s+ s-1})}\\
              &=q^{sr}\gausspol{r+s}{s}_{q^{-1}}
       \end{split}
       \end{equation}
       Of course, each one of the subspaces $N$ of dimension $s$, has codimension $r$ in $N$, i.e. the quotient $M/N$ is of type $(1^r)$. In conclusion, we get
       \begin{equation}
           G_{(1^r)(1^s)}^{(1^{r+s})}=q^{sr}\gausspol{r+s}{s}_{q^{-1}}.
       \end{equation}
\begin{dang}
    \begin{example}
    $G_{(r)(s)}^{(r+s)}$ counts the number of submodules $N$ of type $(s)$ and cotype $(r)$ of $M=R/(p^{(r+s)})$. We claim that 
    \begin{equation}
        G_{(r)(s)}^{(r+s)}=1
    \end{equation}
    \end{example}
\end{dang}
From the isomorphism theorem, each submodule is of the form $N=(p^m)/(p^{r+s})$ for some $m\in\{0,1,\dots,r+s\}$. But, $N$ is of type $(s)$ if and only if $m=r$. Therefore $N=(p^r)/(p^{r+s})$. We must check that $N$ has cotype $(r)$ in $M$. But again, by the isomorphism theorem, $M/N\cong R/(p^{r})$ is of type $(r)$. Therefore the unique possibility is $N=(p^r)/(p^{r+s})$ and we conclude that $G_{(r)(s)}^{(r+s)}=1$.\par
 From the previous examples we get that $G_{(1)(1)}^{(2)}=1$ and 
 \begin{equation}
     G_{(1)(1)}^{(1^2)}=q\smallgausspol{2}{1}_{q^{-1}}=q(1-q^{-2})/(1-q^{-1})=q(1+q^{-1})=q+1
 \end{equation}
 Therefore, we can compute the product
        \begin{equation}
            u_{(1)}u_{(1)}=u_{(2)}+(q+1)u_{(1^2)}.
        \end{equation}
\begin{dang}
    \begin{definition}\cite{macdonald1998symmetric}
        Let $M$ be a fixed $R-$module of type $\lambda$ and partitions $\mu_i,\;i=1,\dots n$. We will write $G_{\mu_1\dots\mu_n}^\lambda$ to denote the number of chains
        \begin{equation}
            (0)=M_n\subset M_{n-1}\subset \cdots\subset M_1\subset M_0=M
        \end{equation}
        such that $M_{i-1}/M_{i}$ is of type $\mu_i$ for each $i=1,\dots,n$.
    \end{definition}
\end{dang}
\begin{dang}
    \begin{lemma}\label{expression_for_the_associativity_of_hall_algebra}
        We have
        \begin{equation}
            \sum_{\lambda}G_{\mu_1\dots\mu_r}^\lambda G_{\lambda\mu_{r+1}\dots\mu_{r+s}}^\xi=G_{\mu_1\dots\mu_{r+s}}^\xi= \sum_{\lambda}G_{\mu_1\dots\mu_{r-1}\lambda}^{\xi} G_{\mu_{r}\dots\mu_{r+s}}^\lambda
        \end{equation}
    \end{lemma}
\end{dang}
Proof: First, let's focus on the left equality. We will compute $G_{\mu_1\dots\mu_{r+s}}^\xi$ by dividing in cases according to the type of $M/M_r$. Clearly, each case will produce a different chain of submodules. Now, fix $\lambda$. By definition, there are $G_{\lambda\mu_{r+1}\dots\mu_{r+s}}^\xi$ ways to construct a chain of the form
\begin{equation}\label{oldchain_in_lemma_about_associativity_of_hall_algebra}
    (0)=M_{r+s}\subset \cdots\subset M_r\subset M
\end{equation}
such that $M_{i-1}/M_i$ is of type $\mu_i$ for each $i=r+1,\dots,r+s$ while $M/M_r$ is of type $\lambda$. Also, there are $G_{\mu_1\dots\mu_r}^\lambda$ ways to construct a chain
\begin{equation}\label{chain_in_lemma_about_associativity_of_hall_algebra}
    (0) = M_r^\prime \subset\cdots\subset M_1^\prime\subset M_0^\prime=M/M_r
\end{equation}
where $M_{i-1}^\prime/M_i^\prime$ is of type $\mu_i$ for each $i=1,\dots,r$. But, by the isomorphism theorem, for each $i$, there is a unique submodule $M_i\subset M$ such that $M/M_i\cong  M_i^\prime$. So, for each chain~(\ref{chain_in_lemma_about_associativity_of_hall_algebra}) there is a unique chain of submodules
\begin{equation}\label{newchain_in_lemma_about_associativity_of_hall_algebra}
    M_r\subset \dots M_1\subset M_0=M
\end{equation}
such that $M_{i-1}/M_i\cong M_{i-1}^\prime/M_i^\prime$ is of type $\mu_i$ for each $i=1,\dots,r$, while $M/M_r$ is of type $\lambda$. Conversely, each chain~(\ref{newchain_in_lemma_about_associativity_of_hall_algebra}) give a chain of the form~(\ref{chain_in_lemma_about_associativity_of_hall_algebra}) by taking the quotient by $M_r$. Therefore, there exists $G_{\mu_1\dots\mu_r}^\lambda$ ways to complete the sequence~(\ref{oldchain_in_lemma_about_associativity_of_hall_algebra}) and we conclude that there exists $G_{\mu_1\dots\mu_r}^\lambda G_{\mu_1\dots\mu_{r+s}}^\xi$ ways to construct a chain $(0)=M_{r+s}\subset\cdots\subset M_1\subset M_0=M$ such that $M_{i-1}/M_i$ is of type $\mu_i$ for each $i=1,\dots,r+s$ while $M/M_r$ is of type $\lambda$. Summing these numbers for all $\lambda$, we get $G_{\mu_1\dots\mu_{r+s}}^\xi$. The same reasoning can be applied to obtain the right equation of the lemma.
\begin{dang}
    \begin{theorem}\cite{macdonald1998symmetric}
          The Hall algebra $H(R)$ is associative.
    \end{theorem}
\end{dang}
Proof: This follows from the previous lemma for $r=2$ and $s=1$.
\begin{dang}
    \begin{theorem}\cite{macdonald1998symmetric}
        The Hall algebra $H(R)$ is commutative.
    \end{theorem}
\end{dang}
Proof: See~(\cite{macdonald1998symmetric}, chapter II).


\begin{dang}
    \begin{definition}
        Let $n\geq1$ be a positive integer and $\mathcal{P}_n$ be the set of partitions of $n$. The \textit{reverse lexicographic ordering} $L_n$ is the subset of $\mathcal{P}_n\times\mathcal{P}_n$ consisting of pairs $(\mu,\lambda)$ such that the first nonvanishing entry $\mu_i-\lambda_i$ is positive. We write $\mu>\lambda$ if $(\mu,\lambda)\in L_n$.
    \end{definition}
\end{dang}
For example, for $n=5$ we have
\begin{equation}
    (5)>(4,1)>(3,2)>(3,1^2)>(2^2,1)>(2,1^3)>(1^5)
\end{equation}
Using this, we put a total order on $\mathcal{P}$ by $\mu>\lambda$ if $\mu\in\mathcal{P}_n$ and $\lambda\in\mathcal{P}_m$ with $n>m$ or if $(\mu,\lambda)\in L_n$ in the case $\mu,\lambda\in\mathcal{P}_n$. With this total order, we are able to work with matrices $(A_{\mu\lambda})$ indexed by $\mathcal{P}$. In this case, multiplication is defined as usual $(AB)_{\mu\lambda}=\sum_\nu A_{\mu\nu}B_{\nu\lambda}$.
\begin{dang}
    \begin{theorem}\label{free generators of the hall algebra}\cite{macdonald1998symmetric}
        Let $\lambda=(\lambda_1,\dots,\lambda_r)$ be a partition and $v_{\lambda^\prime}\bydef v_{\lambda_1}\cdots v_{\lambda_r}$ where $v_i\bydef u_{(1^i)}$. Then, the set $\{v_{\lambda}|\;\lambda\in\mathcal{P}\}$ forms a $\bb{Z}-$basis of $H(R)$.
    \end{theorem}
\end{dang}
Proof: Fix a partition $\lambda$ and let $\lambda^\prime = (\lambda_1^\prime,\dots,\lambda_s^\prime)$ be the conjugate of $\lambda$. By the lemma~(\ref{expression_for_the_associativity_of_hall_algebra}) we get that 
\begin{equation}\label{vlambda_to_ulambda}
    v_{\lambda^\prime}\bydef v_{\lambda_1^\prime}\cdots v_{\lambda_r^\prime}=\sum_{\xi}G_{(1^{\lambda_1^\prime})\dots(1^{\lambda_s^\prime})}^{\;\xi}u_{\xi}
\end{equation}
To ease notation, let's write $a_{\lambda\xi}\bydef G_{(1^{\lambda_1^\prime})\dots(1^{\lambda_s^\prime})}^{\;\xi}$. Now, fix a module $M$ of type $\xi$. By definition, $a_{\lambda\xi}$ is the number of chains
\begin{equation}
    (0)=M_s\subset M_{s-1}\cdots\subset M_1\subset M_0=M
\end{equation}
such that $M_{i-1}/M_i$ is of type $(1^{\lambda_i^\prime})$ for each $i=1,\dots,s$. Suppose that such a sequence exist. Since each $M_{i-1}/M_i$ is elementary, we have $p(M_{i-1}/M_i)=(0)\iff pM_{i-1}\subset M_i$ for each $i=1,\dots,s$. Therefore $p^iM\subset M_i$ and $l(p^iM)\leq l(M_i)$. So
\begin{equation}
    \begin{split}
        &l(M/p^iM)\geq l(M/M_i)\iff \\
        &l(M)-l(p^iM)\geq l(M)-l(M_i) \iff\\
        &\sum_{j=1}^{i}l(p^{j-1}M/p^jM)\geq \sum_{j=1}^{i}l(M_{j-1}/M_j)\iff\\
        &\sum_{j=1}^{i}\xi_j^\prime\geq \sum_{j=1}^{i}\lambda_j^\prime\quad\text{ (by theorem~(\ref{dim_of_M_over_pM})}) 
    \end{split}
\end{equation}
and the fact that $l(M_{j-1}/M_j)=|(1^{\lambda_j^\prime})|=\lambda_j^\prime$. This is true for each $i=1,\dots, s$. Therefore $\xi^\prime\geq\lambda^\prime\iff \xi\leq \lambda$ which implies that $a_{\lambda\xi}\neq0\iff \xi\leq \lambda$. Conversely, if $\xi=\lambda$, we see by the previous equation that $l(M/p^iM)= l(M/M_i)$, which happens only if $l(p^iM)=l(M_i)\Rightarrow p^iM=M$. So in this case, there is only one decreasing chain $(M_i)$ satisfying that $M_{i-1}/M_i$ is of type $(1^{\lambda_i^\prime})$, which is the chain $(p^iM)$. This shows that $a_{\lambda\lambda}=1$. So the change of basis matrix $(a_{\lambda\xi})$ is unitriangular, and therefore unimodular, which implies that we can solve equation~(\ref{vlambda_to_ulambda}) to express the $u_\xi$ in term of the $v_\lambda$. This proves the theorem.\par
One of the consequences of this theorem is that the structure constants $G_{\mu\nu}^\lambda$ are completely determined by the structure constants of the form $G_{\mu(1^m)}^\lambda$. So the next task will be to compute these structure constants explicitly.
\begin{dang}
    \begin{lemma}\label{cotype_contained_in_type}\cite{macdonald1998symmetric}
        Let $M$ be a finite $R-$module of type $\lambda$ and $N\subset M$ a submodule of cotype $\mu$ in $M$. Then $\mu\subset \lambda$.
    \end{lemma}
\end{dang}
Proof: We need to show that $\mu_i\leq\lambda_i$ for all $i\geq1$. Since $p^i(M/N)$ is a submodule of $M/N$, by the isomorphism theorem, there exists a submodule $L_i\subset M$ such that $M/L_i=p^i(M/N)$. In this case $L^i=\pi^{-1}(p^i(M/N))=p^iM+N$ where $\pi$ in the natural projection of $M$ onto $M/N$. Therefore
\begin{equation}
    \begin{split}
        \frac{p^{i-1}(M/N)}{p^i(M/N)}&=\frac{(p^{i-1}M+N)/N}{(p^{i}M+N)/N}\\
        &\cong\frac{p^{i-1}M+N}{p^{i}M+N}\\
        &=\frac{p^{i-1}M+p^iM+N}{p^{i}M+N}\quad\text{ (because } p^iM\subset p^{i-1}M)\\
        &\cong\frac{p^{i-1}M}{p^{i-1}M\cap(p^{i}M+N)}
    \end{split}
\end{equation}
Therefore
\begin{equation}
    l\left( \frac{p^{i-1}(M/N)}{p^i(M/N)}\right) = l(p^{i-1}M)-l(p^{i-1}M\cap(p^{i}M+N))
\end{equation}
But $p^{i}M\subset p^{i-1}M\cap(p^{i}M+N)$, so 
\begin{equation}
     \mu_i^\prime= l\left( \frac{p^{i-1}(M/N)}{p^i(M/N)}\right) \leq l(p^{i-1}M)-l(p^{i}M)=l\left(\frac{p^{i-1}M}{p^{i}M}\right)=\lambda_i^\prime
\end{equation}
and we conclude that $\mu_i\leq\lambda_i$ for each $i\geq1\Rightarrow \mu\subset\lambda$.
\begin{dang}
    \begin{definition}\cite{macdonald1998symmetric}
        Let $M$ be a $R-$module. The set $S=\{x\in M|\;px=0\}$ is called the \textit{socle} of $M$.
    \end{definition}
\end{dang}
It's clear by the definition that $S\subset M$ is the unique maximal elementary submodule of $M$. Now, if $M$ is of type $\lambda=(\lambda_1,\dots,\lambda_r)$, i.e.,
\begin{equation}
    M=\bigoplus_{i=1}^r\frac{R}{(p^{\lambda_i})}
\end{equation}
then $S=\bigoplus_{i=1}^r(p^{\lambda_i-1})/(p^{\lambda_i})$ and $M/S\cong\bigoplus_{i=1}^r(R/(p^{\lambda_i-1}))$. So $M/S$ is of type $\Tilde{\lambda}=(\lambda_1-1,\dots,\lambda_r-1)$.
\begin{dang}
    \begin{lemma}\label{alpha subset beta subset lambda}\cite{macdonald1998symmetric}
        Let $M$ be a finite $R-$module of type $\lambda$ and $L\subset N\subset M$ where $L$ (resp. $N$) is of cotype $\beta$ (resp. $\alpha$) in $M$. With these conditions $\alpha\subset \beta\subset \lambda$.
    \end{lemma}
\end{dang}
Proof: We have the isomorphism $M/N\cong (M/L)/(N/L)$. So $N/L$ is a submodule of cotype $\alpha$ in $M/L$ . But $M/L$ is of type $\beta$, which implies that $\alpha\subset\beta$ by lemma~(\ref{cotype_contained_in_type}). Moreover, the same lemma implies that $\beta\subset\lambda$ and the result follows.
\begin{dang}
    \begin{lemma}\label{vertical strip lemma}\cite{macdonald1998symmetric}
        Let $M$ be a finite $R-$module of type $\lambda$ and $N\subset M$ be an elementary submodule of cotype $\mu$. Then $\lambda-\mu$ is a \textit{vertical strip}, i.e. $\lambda_i-\mu_i\in\{0,1\}$ for each $i\geq1$. 
    \end{lemma}
\end{dang}
Proof: Since $N$ is elementary, we have $N\subset S$. Therefore, from the previous lemma $\Tilde{\lambda}\subset\mu\subset\lambda\Rightarrow \lambda_i-1\leq\mu_i\leq\lambda_i$ for each $i\geq1$. This shows that $\lambda_i-\mu_i\in\{0,1\}$ for each $i\geq1$.

\begin{dang}
    \begin{theorem}\cite{macdonald1998symmetric}
        Let $M$ be a finite $R-$module of type $\lambda$ and $N\subset M$ a submodule of cotype $\alpha$ in $M$. Let $\beta$ be a partition such that $\alpha\subset \beta\subset\lambda$. The number of submodules $L\subset N$ of cotype $\beta$ in $M$ is given by
        \begin{equation}
            h_{\alpha\beta\lambda}(q) \bydef q^{d(\alpha,\beta,\lambda)}\prod_{i\geq1}\gausspol{\lambda_i^\prime-\alpha_i^\prime}{\lambda_i^\prime-\beta_i^\prime}_{q^{-1}}
        \end{equation}
        where $d(\alpha,\beta,\lambda)=\sum_{j\geq i}(\beta_j^\prime-\alpha_j^\prime)(\lambda_i^\prime-\beta_i^\prime)$.
    \end{theorem}
\end{dang}
Proof: Let $L\subset N$ be a submodule (not necessarily of cotype $\beta$ in $M$). Also, define $N_i\bydef N\cap p^iM,\;L_i\bydef L\cap p^iM\;(= L\cap N_i),\; n_i\bydef l(N_i)$ and $l_i\bydef l(L_i)$. Then
\begin{equation}
    \begin{split}
        l(L_{i-1}/L_i)&=l(L/L_i)-l(L/L_{i-1})\\
        &=l(L/(L\cap p^iM))-l(L/L\cap p^{i-1}M)\\
        &=l\left(\frac{L+p^iM}{p^iM}\right)-l\left(\frac{L+p^{i-1}M}{p^{i-1}M}\right)\\
        &=l(p^{i-1}M/p^iM)-l((L+p^{i-1}M)/(L+p^iM))\\
        &=\lambda_i^\prime - l\left(\frac{(L+p^{i-1}M)/L}{(L+p^{i}M)/L}\right)\\
        &=\lambda_i^\prime - l\left(\frac{p^{i-1}(M/L)}{p^i(M/L)}\right).
    \end{split}
\end{equation}
Now, let $\theta\bydef\lambda-\beta$. The previous equation shows that $L$ is of cotype $\beta$ in $M$ if and only if $l(L_{i-1}/L_i)=\lambda_i^\prime-\beta_i^\prime=\theta_i^\prime$. So we see that for each  submodule $L$ of cotype $\beta$ in $M$, there is a sequence $(L_i)_{i\geq0}$ (with $L=L_0$) which satisfies the following conditions:
\begin{equation}
    \textbf{(1) } L_i=L_{i-1}\cap N_i\quad\text{ and }\quad\textbf{(2) }l(L_{i-1}/L_i)=\theta_i^\prime\quad\forall i
\end{equation}
Conversely, for each sequence $(L_i)_{i\geq0}$ satisfying (1) and (2), we have $L_i = L_0\cap N_i$ while $l(L_{i-1}/L_i)=\theta_i^\prime$. So $L_0$ is a submodule of cotype $\beta$ in $M$. This shows that, the map $(L_i)_{i\geq0}\mapsto L_0$ gives a one-to-one correspondence between sequences satisfying (1) and (2) and submodules of cotype $\beta$ in $M$. Therefore, we will focus on enumerating such sequences.
\par
The strategy will be the following: For a given $i\geq1$ and $L_i\subset N_i$, we will count the number of submodules $L_{i-1}\subset N_{i-1}$ such that (1) and (2) are satisfied for $i$. We will see that $L_{i-1}$ is generated by $L_i$ and some linearly independent set modulo $N_i$ in $N_{i-1}$, with $\theta_i^\prime$ elements.\par
Now we need to make an observation: Let $L_j\subset N_j\;(j=i-1,i)$ be two submodules satisfying (1) and (2) for $i$. Since $L_{i-1}/L_i$ is a vector space, by the property (1) and (2), we get that $\text{dim}_{\bb{F}}(L_{i-1}/L_i)=l(L_{i-1}/L_i)=\theta_i^\prime$. Similarly, $N_{i-1}/N_i$ is a vector space of dimension $l(N_{i-1}/N_i)=n_{i-1}-n_i$. But, since $N_i=N\cap p^iM$, we get
\begin{equation}
\begin{split}
        p^iM/N_i &\cong (N+p^iM)/N\cong p^i(M/N) \\
        \Rightarrow l(N_i)&=l(p^iM)-l(p^i(M/N))\\
        &=\sum_{j>i}l(p^{j-1}M/p^jM)-l(p^{j-1}(M/N)/p^j(M/N))\\
        &=\sum_{j>i}(\lambda_j^\prime-\alpha_j^\prime)\quad\text{ (because of the lemma~(\ref{dim_of_M_over_pM})})
\end{split}
\end{equation}
The last equation implies that $n_{i-1}-n_i=\lambda_i^\prime-\alpha_i^\prime$. To ease notation, define $\varphi\bydef \beta-\alpha$, so we can write $\theta_i^\prime+\varphi_i^\prime=\lambda_i^\prime-\alpha_i^\prime$ for the dimension of $L_{i-1}/L_i$.\par
With these observations, we proceed to count the submodules $L_{i-1}$. By Lemma~(\ref{number o LI sequence modulo N}), there exists
\begin{equation}
   A_i=(q^{\theta_i^\prime+\varphi_i^\prime+n_{i}}-q^{n_{i}})(q^{\theta_i^\prime+\varphi_i^\prime+n_{i}}-q^{n_{i}+1})\dots(q^{\theta_i^\prime+\varphi_i^\prime+n_{i}}-q^{n_{i}+\theta_i^\prime-1})
\end{equation}
sequences $x=(x_1,\dots,x_{\theta_i^\prime})$ in $N_{i-1}$, which are linearly independent modulo $N_i$. Writing $Rx$ for the $R-$span of $x$, we claim that $L_{i-1}\bydef Rx+L_i$ satisfies properties (1) and (2). \par
First, lets verify the first property. It's clear that $L_i\subset (Rx+L_i)\cap N_i$. Now, if $y\in (Rx+L_i)\cap N_i$, then $y=z+l\in N_i$, where $z\in Rx$ and $l\in N_i$. But this implies that $z\in N_i$, because $L_i\subset N_i$. But $x$ is linearly independent modulo $N_i$. Therefore $z=0$ and $y=l\in N_i$. This shows that $L_i=(Rx+L_i)\cap N_i=L_{i-1}\cap N_i$. For the second property, we just need to see that $L_{i-1}/L_i=(Rx+L_i)/L_i$ is a vector space of dimension $\theta_i^\prime$, because $(x_1+N_i,\dots, x_{\theta_i^\prime}+N_i)$ is a basis of this space.\par
Conversely, if $L_{i-1}\subset L_i$ is such that (1) and (2) is true for $i$, then $L_{i-1}/L_i$ is a vector space of dimension $\theta_i^\prime$. Let $(x_1+L_i,\dots, x_{\theta_i^\prime}+L_i)$ be a basis of $L_{i-1}/L_i$, so $L_{i-1}=Rx+L_i$ where $(x_1,\dots,x_{\theta_i^\prime})$ is a sequence in $L_{i-1}\subset N_{i-1}$ which is linearly independent modulo $L_i$. By the lemma~(\ref{number o LI sequence modulo N}), there are 
\begin{equation}
      B_i=(q^{\theta_i^\prime+l_{i}}-q^{l_{i}})(q^{\theta_i^\prime+l_{i}}-q^{l_{i}+1})\dots(q^{\theta_i^\prime+l_{i}}-q^{l_{i}+\theta_i^\prime-1})
\end{equation}
such sequences. This shows that, for each one of the $A_i$ ways to construct $L_{i-1}$, there are exactly $B_i$ that produces the same module. Therefore, for a fixed $i\geq1$ and $L_i\subset N_i$, the number of modules $L_{i-1}\subset N_{i-1}$ such that $(1)$ and $(2)$ hold for $i$ is
\begin{equation}\label{A/B}
\begin{split}
     \frac{A_i}{B_i}&=\frac{(q^{\theta_i^\prime+\varphi_i^\prime+n_{i}}-q^{n_{i}})(q^{\theta_i^\prime+\varphi_i^\prime+n_{i}}-q^{n_{i}+1})\dots(q^{\theta_i^\prime+\varphi_i^\prime+n_{i}}-q^{n_{i}+\theta_i^\prime-1})}{(q^{\theta_i^\prime+l_{i}}-q^{l_{i}})(q^{\theta_i^\prime+l_{i}}-q^{l_{i}+1})\dots(q^{\theta_i^\prime+l_{i}}-q^{l_{i}+\theta_i^\prime-1})}\\
     &=\frac{(1-q^{-(\theta_i^\prime+\varphi_i^\prime)})(1-q^{-(\theta_i^\prime+\varphi_i^\prime)+1})\dots(1-q^{-(\theta_i^\prime+\varphi_i^\prime)+\theta_i^\prime-1})}{(1-q^{-\theta_i^\prime})(1-q^{-\theta_i^\prime+1})\dots(1-q^{-\theta_i^\prime+\theta_i^\prime-1})}\frac{q^{\theta_i^\prime n_{i-1}}}{q^{\theta_i^\prime l_{i-1}}}\\
     &=q^{\theta_i^\prime(n_{i-1}-l_{i-1})}\gausspol{\theta_i^\prime+\varphi_i^\prime}{\theta_i^\prime}_{q^{-1}}\\
     &=q^{\theta_i^\prime\sum_{j\geq i}\varphi_j^\prime}\gausspol{\theta_i^\prime+\varphi_i^\prime}{\theta_i^\prime}_{q^{-1}}\;\;\text{ (because }\;l_{i-1}=\sum_{j\geq i}\theta_i^\prime,\;\;\text{ which follows from } (2))
\end{split}
\end{equation}
Therefore, the number of sequences $(L_i)_{i\geq0}$, satisfying conditions (1) and (2), is
\begin{equation}
\begin{split}
       &\prod_{i\geq1}q^{\theta_i^\prime\sum_{j\geq i}\varphi_j^\prime}\gausspol{\theta_i^\prime+\varphi_i^\prime}{\theta_i^\prime}_{q^{-1}}\\
       &=q^{\sum_{i\geq1}\theta_i^\prime\sum_{j\geq i}\varphi_j^\prime}\prod_{i\geq1}\gausspol{\theta_i^\prime+\varphi_i^\prime}{\theta_i^\prime}_{q^{-1}}\\
       &=q^{\sum_{j\geq i}\theta_i^\prime\varphi_j^\prime}\prod_{i\geq1}\gausspol{\theta_i^\prime+\varphi_i^\prime}{\theta_i^\prime}_{q^{-1}}\\
       &=q^{d(\alpha,\beta,\lambda)}\prod_{i\geq1}\gausspol{\lambda_i^\prime-\alpha_i^\prime}{\lambda_i^\prime-\beta_i^\prime}_{q^{-1}}\text{ (because }\theta_i^\prime=\lambda_i^\prime-\beta_i^\prime\text{ and }\varphi_i^\prime=\beta_i^\prime-\alpha_i^\prime).
\end{split}
\end{equation}
This completes the proof of the theorem.\par


Now, we will use this theorem to compute structure constants of the form $G_{\beta (1^m)}^\lambda$. To do this, 
we need to make some observations about the previous theorem. First, the condition $\alpha\subset\beta$ is necessary, 
because of the lemma~(\ref{alpha subset beta subset lambda}). Also, if we take $N=S$, the socle of $M$, 
we get that $L\subset S$ if and only if $L$ is elementary. Then $L\subset S$ implies that $L$ is elementary and 
it is of type $(1^m)$ 
for some $m\geq1$. In this case, $L$ exists only if $|\lambda|=|\beta|+|(1^m)|\iff m=|\lambda|-|\beta|$. Also, since $S$ is the socle of $M$, we get that $\alpha=\Tilde{\lambda}=(\lambda_1-1,\lambda_2-1,\dots)$. But also 
\begin{equation}
    \begin{split}
        \Tilde{\lambda}_i^\prime &= \text{Max}\{k\;|\;\Tilde{\lambda}_k\geq i\}\\
          &=\text{Max}\{k\;|\;\lambda_k-1\geq i\}\\
        &=\text{Max}\{k\;|\;\lambda_k\geq i+1\}\\
        &=\lambda_{i+1}^\prime
    \end{split}
\end{equation}
So we have
\begin{equation}
    \prod_{i\geq1}\gausspol{\lambda_i^\prime-\alpha_i^\prime}{\lambda_i^\prime-\beta_i^\prime}_{q^{-1}}=\;\;\; \prod_{i\geq1}\gausspol{\lambda_i^\prime-\lambda_{i+1}^\prime}{\lambda_i^\prime-\beta_i^\prime}_{q^{-1}}
\end{equation}
Now, it remains to compute $d(\alpha,\beta,\lambda)$. Recall the notation $\theta=\lambda-\beta$ and $\varphi=\beta-\alpha=\beta-\Tilde{\lambda}$. So $\varphi_i=\beta_i-\lambda_i+1=1-\theta_i$. Now, we need to calculate the sum $d(\alpha,\beta,\lambda)=\sum_{j\geq i}\varphi_j^\prime\theta_i^\prime=\sum_{j\geq i}\theta_j\varphi_i$. We have
\begin{equation}
    \begin{split}
        \sum_{j\geq i}\theta_j\varphi_i &= \sum_{j\geq i}\theta_j(1-\theta_i)\\
        &=\sum_{j>i}\theta_j(1-\theta_i)\quad(\text{because }\theta_i\in\{0,1\}\text{ by lemma~(\ref{vertical strip lemma})})\\
        &=\sum_{j>i}\theta_j-\sum_{j>i}\theta_j\theta_i\\
        &=\sum_{i\geq1}(i-1)\theta_i - \frac{(\sum_i\theta_i)^2-\sum_i\theta_i^2}{2}\\
        &=n(\lambda)-n(\beta)-\frac{m^2-m}{2}\quad(\text{ because } m =|\theta| \text{ and } \theta_i^2=\theta_i)\\
        &=n(\lambda)-n(\beta)-n((1^m)).
    \end{split}
\end{equation}
So we get the next result.
\begin{dang}
    \begin{theorem}\label{fundamental structure constants theorem}\cite{macdonald1998symmetric}
        Let $M$ be a finite $R-$module of type $\lambda$. Then, for any elementary submodule $L\subset M$ of cotype $\beta$, the skew diagram $\lambda-\beta$ is a vertical strip and the number of such modules which are of type $(1^m)$ is
        \[
        G_{\beta(1^m)}^\lambda=q^{n(\lambda)-n(\beta)-n((1^m))}\prod_{i\geq1}\gausspol{\lambda_i^\prime-\lambda_{i+1}^\prime}{\lambda_i^\prime-\beta_i^\prime}_{q^{-1}}.
        \]
    \end{theorem}
\end{dang}
Theorem~(\ref{fundamental structure constants theorem}) implies that the structure constants $G_{\beta(1^m)}^\lambda$ are determined by the following ``universal'' polynomial
\begin{equation}
    f_{\beta(1^m)}^\lambda(t) \bydef t^{n(\lambda)-n(\beta)-n((1^m))}\prod_{i\geq1}\gausspol{\lambda_i^\prime-\lambda_{i+1}^\prime}{\lambda_i^\prime-\beta_i^\prime}_{t^{-1}}.
\end{equation}
That is, $G_{\beta(1^m)}^\lambda= f_{\beta(1^m)}^\lambda(q)$ for any ring $R$ with residue field of cardinality $q$. This, and theorem~(\ref{free generators of the hall algebra}) show that the structure constants depend only on the cardinality of the residue field of $R$. Also, it is possible to explicitly determine all structure constants $G_{\mu\nu}^\lambda$ (see~\cite{macdonald1998symmetric} chapter II).\par
Theorem~(\ref{fundamental structure constants theorem}) will be useful to show that the space of characters of $GL_n(\fq)$ is a certain infinite tensor product of Hall algebras, which we do in the next chapter.
